% concatenated from icml_2025.tex

\documentclass{article} % For LaTeX2e
\usepackage[accepted]{icml2025} 
% Recommended, but optional, packages for figures and better typesetting:
\usepackage{microtype}
\usepackage{graphicx}
\usepackage{subfigure}
\usepackage{booktabs} % for professional tables
\usepackage{array}
\usepackage{color}
\usepackage{subcaption}
\usepackage{wrapfig}

% hyperref makes hyperlinks in the resulting PDF.
% If your build breaks (sometimes temporarily if a hyperlink spans a page)
% please comment out the following usepackage line and replace
% \usepackage{icml2025} with \usepackage[nohyperref]{icml2025} above.
\usepackage[colorlinks=true, linkcolor=blue, citecolor=blue, urlcolor=blue]{hyperref}


% Attempt to make hyperref and algorithmic work together better:
\newcommand{\theHalgorithm}{\arabic{algorithm}}


\usepackage[utf8]{inputenc} % allow utf-8 input
\usepackage[T1]{fontenc}    % use 8-bit T1 fonts
\usepackage{hyperref}       % hyperlinks
\usepackage{url}            % simple URL typesetting
\usepackage{booktabs}       % professional-quality tables
\usepackage{amsfonts}       % blackboard math symbols
\usepackage{nicefrac}       % compact symbols for 1/2, etc.
\usepackage{microtype}      % microtypography
\usepackage{xcolor}         % colors

\usepackage{algorithm}
% \usepackage{algpseudocode}

% For theorems and such
\usepackage{amsmath}
\usepackage{amssymb}
\usepackage{mathtools}
\usepackage{amsthm}
\usepackage{bm}

\theoremstyle{plain}
\newtheorem{theorem}{Theorem}[section]
\newtheorem{proposition}[theorem]{Proposition}
\newtheorem{Lemma}[theorem]{Lemma}
\newtheorem{corollary}[theorem]{Corollary}
\theoremstyle{definition}
\newtheorem{definition}[theorem]{Definition}
\newtheorem{assumption}[theorem]{Assumption}
\theoremstyle{remark}
\newtheorem{remark}[theorem]{Remark}
\newcommand{\mb}[1]{\mathbf{#1}}
\newcommand{\red}[1]{{\color{red} #1}}

\bibliographystyle{abbrvnat}

\begin{document}

\twocolumn[
\icmltitle{Revisiting Convergence: Shuffling Complexity Beyond Lipschitz Smoothness}

\begin{icmlauthorlist}
\icmlauthor{Qi He}{yyy}
\icmlauthor{Peiran Yu}{zzz}
\icmlauthor{Ziyi Chen}{yyy}
\icmlauthor{Heng Huang}{yyy}
\end{icmlauthorlist}

\icmlaffiliation{yyy}{Department of Computer Science, University of Maryland, College Park.}
\icmlaffiliation{zzz}{Department of Computer Science and Engineering, University of Texas Arlington}

\icmlcorrespondingauthor{Qi He}{qhe123@umd.edu}
\icmlcorrespondingauthor{Heng Huang}{heng@umd.edu}

\vskip 0.3in
]

\printAffiliationsAndNotice{}

\begin{abstract}
  Shuffling-type gradient methods are favored in practice for their simplicity and rapid empirical performance. Despite extensive development of convergence guarantees under various assumptions in recent years, most require the Lipschitz smoothness condition, which is often not met in common machine learning models. We highlight this issue with specific counterexamples. To address this gap, we revisit the convergence rates of shuffling-type gradient methods without assuming Lipschitz smoothness. Using our stepsize strategy, the shuffling-type gradient algorithm not only converges under weaker assumptions but also match the current best-known convergence rates, thereby broadening its applicability. We prove the convergence rates for nonconvex, strongly convex, and non-strongly convex cases, each under both random reshuffling and arbitrary shuffling schemes, under a general bounded variance condition. Numerical experiments further validate the performance of our shuffling-type gradient algorithm, underscoring its practical efficacy.
\end{abstract}
\section{Introduction}

Gradient-based optimization has always been a critical area due to its extensive practical applications in machine learning, including reinforcement learning \citep{sutton2018reinforcement}, hyperparameter optimization \citep{feurer2019hyperparameter}, and large language models \citep{radford2018improving}. While numerous gradient-based algorithms have been developed for convex functions \citep{nemirovskij1983problem, nesterov2013gradient, d2021acceleration}, research on nonconvex functions has become particularly active in recent years, driven by advances in deep learning. Notably, with unbiased stochastic gradients and bounded variance, SGD achieves an optimal complexity of $\mathcal{O}(\epsilon^{-4})$ \citep{ghadimi2013stochastic}, which matches the lower bound established by \citet{arjevani2023lower}.

In practice, however, random shuffling-type methods have demonstrated advantages over traditional SGD. 
These methods, which involve iterating over data in a permuted order rather than drawing i.i.d. samples, 
are widely used in deep learning frameworks due to their simplicity and computational efficiency. 
Empirical studies have shown that such methods often outperform standard SGD in terms of convergence speed and 
generalization ability \citep{Bo09, Bo12}. 
Notably, shuffling-based training has been found to mitigate gradient variance, 
leading to smoother optimization trajectories compared to purely stochastic updates \citep{GuOzPa21, HaSu19}. 

Theoretical analyses of shuffling-type methods have gained significant attention in recent years, 
aiming to explain their empirical success while tackling the unique challenges posed by 
the lack of independence between consecutive updates. Many early works focused on convex optimization, 
often relying on strong convexity assumptions to establish linear or sublinear convergence rates 
\citep{GuOzPa21, HaSu19, SaSh20}. 
More recently, research efforts have extended beyond the convex setting \citep{mishchenko2020random, nguyen2021unified, koloskova2023shuffle},
suggesting that random reshuffling can exhibit faster convergence than SGD even in nonconvex settings 
under certain conditions, such as when gradient noise is structured or exhibits low variance. 

Despite these advances, most theoretical analyses rely on the Lipschitz smoothness assumption, 
which imposes restrictive upper and lower bounds on the gradient variations. 
While this assumption holds in many standard optimization settings, it fails to capture 
a broad class of important machine learning models, including deep language models \citep{zhang2019gradient}, 
phase retrieval \citep{chen2023generalized}, and distributionally robust optimization \citep{chen2023generalized}. 
As a result, many practical scenarios remain outside the scope of existing theoretical guarantees. 
To address this gap, our work develops new techniques to analyze the convergence of 
shuffling-type gradient algorithms under relaxed smoothness assumptions, 
aiming to provide a more comprehensive theoretical foundation applicable to a wider range of machine learning models.

% In practice, however, random shuffling-type methods have demonstrated superiority over SGD. These methods are not only easier and faster to implement but also show faster convergence rates, as evidenced by experiments cited in \citet{Bo09, Bo12}. Theoretical studies on shuffling-type methods have been conducted in various settings in recent years, presenting unique challenges due to the lack of independence between most neighboring steps. While much of this research assumes strong convexity \citep{GuOzPa21, HaSu19, SaSh20}, studies such as \citet{nguyen2021unified, koloskova2023shuffle, mishchenko2020random} have also explored applications in nonconvex scenarios. 

% Although theoretical analysis has been conducted in many settings of shuffling-type gradient algorithms, \textbf{most of these works require Lipschitz smoothness assumption, which requires restrictive quadratic lower and upper bounds and thus cannot cover many popular machine learning models} such as language model \citep{zhang2019gradient}, phase retrieval  \citep{chen2023generalized}, distributionally robust optimization \citep{chen2023generalized}, etc. We will demonstrate counterexamples in more detail in Section \ref{prelim}. To fill this gap, in this paper, we aim at analyzing the convergence rate of shuffling-type gradient algorithm under relaxed mild smoothness assumptions for both convex and nonconvex cases.

We consider the following finite sum minimization problem:
\begin{equation}
\min_{w \in \mathbb{R}^d} \left\{ F(w) := \frac{1}{n} \sum_{i=1}^n f(w;i) \right\}, \tag{P}
\end{equation}
where $f(\cdot;i) : \mathbb{R}^d \rightarrow \mathbb{R}$ is smooth and possibly nonconvex for $i \in [n] := \{1, \ldots, n\}$. Problem (P) covers empirical loss minimization as a special case, therefore can be viewed as formulation for many machine learning models, such as logistic regression, reinforcement learning, and neural networks.

We summarize our main contributions as follows:

\begin{itemize}
    \item We proved the convergence of the shuffling-type gradient algorithm under non-uniform smoothness assumptions, where the Hessian norm is bounded by a sub-quadratic function $\ell$ of the gradient norm. With specific stepsizes and a general bounded variance condition, we achieved a total complexity of $\mathcal{O}(n^\frac{p+1}{2}\epsilon^{-3})$ gradient evaluations for the nonconvex case with random reshuffling, and $\mathcal{O}(n^{\frac{p}{2}+1}\epsilon^{-3})$ for arbitrary scheme, where $0\leq p<2$ is the degree of function $\ell$ . These results match those with Lipschitz smoothness assumptions in \citet{nguyen2021unified} when $p=0$ and $\ell$-smoothness degenerates to Lipschitz smoothness.
    \item For the strongly convex case, we established a complexity of $\widetilde{\mathcal{O}}(n^\frac{p+1}{2}\epsilon^{-\frac{1}{2}})$ for random reshuffling. In the non-strongly convex case, the complexity is $\mathcal{O}(n^\frac{p+1}{2}\epsilon^{-\frac{3}{2}})$ for random reshuffling. Without assuming bounded variance, we established complexity of $\widetilde{\mathcal{O}}(n\epsilon^{-\frac{1}{2}})$ for arbitrary scheme in strongly convex case, and $\mathcal{O}(n\epsilon^{-\frac{3}{2}})$ in non-strongly convex case.
    \item We conducted numerical experiments to demonstrate that the shuffling-type gradient algorithm converges faster than SGD on two important non-Lipschitz smooth applications.
\end{itemize}

\section{Preliminaries}
\label{prelim}

\subsection{Shuffling-Type Gradient Algorithm}

In practice, the random shuffling method has demonstrated its superiority over SGD, as shown in \citet{Bo09} and \citet{Bo12}. Specifically, \citet{Bo09} shows that shuffling-type methods achieve a convergence rate of approximately $\mathcal{O}(1/T^2)$, where $T$ is the iteration count. Beyond shuffling-type stochastic gradient methods, variants such as SVRG have been applied in various scenarios, including decentralized optimization, as discussed in \citet{Sh16} and \citet{DeGo16}.

The analysis of shuffling-type methods has a long history. For convex cases, \citet{GuOzPa21} demonstrated that when the objective function is a sum of quadratics or smooth functions with a Lipschitz Hessian, and with a diminishing stepsize, the average of the last update in each epoch of RGA converges strictly faster than SGD with probability one. Additionally, they showed that when the number of epochs $T$ is sufficiently large, the Reshuffling Gradient Algorithm (RGA) asymptotically converges at a rate of $\mathcal{O}(1/T^2)$. Similarly, \citet{nguyen2021unified} established a convergence rate of $\mathcal{O}(1/T^2)$ for strongly convex and globally $L$-smooth functions. Furthermore, with uniform sampling and a bounded variance assumption or convexity on each component function, they showed that the convergence rate can be improved to $\mathcal{O}(1/nT^2)$.

In contrast, there is not much research on nonconvex cases. For example, \citet{nguyen2021unified} demonstrated a convergence rate of $\mathcal{O}(T^{-2/3})$; \citet{koloskova2023shuffle} proved a convergence rate of $\mathcal{O}\left(\frac{1}{T}+\min\left\{\left(\frac{n\sigma}{T}\right)^{\frac23},\left(\frac{n\sigma^2}{T}\right)^{\frac12}\right\}\right)$ for single shuffling gradient method; \citet{wang2024provable} analyzed Adam algorithm with random reshuffling scheme under $(L_0,L_1)$-smoothness but they did not explicitly show the convergence rate.

\subsection{Counterexamples}

Although $L$-smoothness is empirically false in LSTM \citep{zhang2019gradient} and Transformers \citep{crawshaw2022robustness}, we give some concrete counterexamples to demonstrate the popularity of non-Lipschitz functions in this section. First we give two machine learning examples, then we mention some common non-Lipschitz functions.

\paragraph{Example 1.}\label{eg:DRO} The first example is distributionally robust optimization (DRO), which is a popular optimization framework for training robust models. DRO is introduced to deal with the distribution shift between training and test datasets. In \citep{levy2020large}, it is formulated equivalently as follows.
\begin{equation}
    \min_{w\in \mathcal{W}, \theta \in \mathbb{R}} L(w,\theta) := \mathbb{E}_{\xi \sim P} \psi^*\left( \frac{\ell(w; \xi) - \theta}{\lambda} \right) + \theta,\label{eq:DRO}
\end{equation}
where $w$ and $\theta$ are the parameters to be optimized, $\xi$ is a sample randomly drawn from data distribution $P$, $\ell(w; \xi)$ is the loss function, $\psi^*$ is the conjugate function of the divergence function $\psi$ we choose to measure the difference between distributions, and $\lambda>0$ is the regularization coefficient. It is proved in \citet{jin2021non} that $L(w,\theta)$ is not always Lipschitz-smooth even if $\ell(w;\xi)$ is Lipschitz-smooth and the variance is bounded.

\paragraph{Example 2.}\label{eg:phase} The second example is the phase retrieval problem. Phase retrieval is a nonconvex problem in X-ray crystallography and diffraction imaging \citep{drenth2007principles, miao1999extending}. The goal is to recover the structure of a molecular object from intensity measurements. Let \( x \in \mathbb{R}^d \) be the true object and \( y_r = |a_r^\top x|^2 \) for \( r = 1, \ldots, m \), where \( a_r \in \mathbb{R}^d \). The problem is to solve:

\begin{align}
\min_{z \in \mathbb{R}^d} f(z) := \frac{1}{2m} \sum_{r=1}^m (y_r - |a_r^\top z|^2)^2.\label{obj_phase}
\end{align}

This objective function is a high-order polynomial in high-dimensional space, thus it does not belong to the \( L \)-smooth function class.

\paragraph{Example 3.} There are many common functions that are not Lipschitz smooth, including polynomial functions with order $>2$, exponential functions, logarithmic functions and rational functions.

\subsection{Relaxation of Lipschitz Smoothness}

Because of the existence of these counterexamples, people have recently been investigating about smoothness assumptions that are more general than the traditional Lipschitz smoothness. In \citet{zhang2019gradient}, $(L_0,L_1)$-smoothness was proposed as the first relaxed smoothness notion motivated by language modeling. It is defined as below:
\begin{definition}
    ($(L_0,L_1)$-smoothness) A real-valued differentiable function $f$ is $(L_0,L_1)$-smooth if there exist constants $L_0, L_1>0$ such that
    $$\|\nabla^2 f(w)\| \leq L_0+L_1\|\nabla f(w)\|.$$
\end{definition}

Lipschitz smoothness can be viewed as a special case of $(L_0,L_1)$ smoothness when $L_1=0$. Under $(L_0,L_1)$-smoothness assumption, various convergence algorithms have been developed including clipped or normalized GD/SGD
\citep{zhang2019gradient}, momentum accelerated clipped GD/SGD \citep{zhang2020improved}, ADAM \citep{wang2022provable} and variance-reduced clipping \citep{reisizadeh2023variance} with optimal sample complexity on stochastic non-convex optimization. 

Other relaxed smoothness assumptions include asymmetric generalized smoothness motivated by distributionally robust optimization \citep{jin2021non} and its extension to $\alpha$-symmetric generalized smoothness \citep{chen2023generalized} and $\ell$-smoothness \citep{li2023convex}. In this paper, we use the definition of $\ell$-smoothness as below:

\begin{definition}
\label{def:gs}
    ($\ell$-smoothness) A real-valued differentiable function $f$ is $\ell$-smooth if there exists some non-decreasing continuous function $\ell:[0,+\infty)\rightarrow (0,+\infty)$ such that for any $w\in \text{dom}(f)$ and constant $C>0$, $\mathcal{B}(w,
    \frac{C}{\ell(\|\nabla f(w)\|+C)})\subseteq \text{dom}(f)$; and for any $w_1,w_2\in \mathcal{B}(w, \frac{C}{\ell(\|\nabla f(w)\|+C)})$,
    $$\|\nabla f(w_1)-\nabla f(w_2)\|\leq \ell (\|\nabla f(w)\|+C)\cdot\|w_1-w_2\|.$$
\end{definition} 

Another equivalent definition of $\ell$-smoothness in \citet{li2023convex} is:

\begin{definition}
\label{def:gs2}
A real-valued differentiable function \( f  \) is \(\ell\)-smooth for some non-decreasing continuous function \( \ell : [0,+\infty) \to (0,+\infty) \) if 
\[
\|\nabla^2 f(x)\| \leq \ell(\|\nabla f(x)\|)
\]
almost everywhere.
\end{definition}

In the proof of sections we focus on Definition \ref{def:gs}, though Definition \ref{def:gs2} provides a clearer perspective on the relationship between \(\ell\)-smoothness and other smoothness notions. Both \((L_0, L_1)\)-smoothness and traditional Lipschitz smoothness can be seen as special cases of \(\ell\)-smoothness. 
It is straightforward to verify that the loss functions in phase retrieval and distributionally robust optimization (DRO) satisfy \(\ell\)-smoothness. 
While extensive research has been conducted based on \((L_0, L_1)\)-smoothness, comparatively less work has been done under the more general \(\ell\)-smooth framework. 
\citet{xiandelving} established convergence results for generalized GDA/SGDA and GDmax/SGDmax in minimax optimization, while \citet{zhang2024convergence} analyzed MGDA for multi-objective optimization under \(\ell\)-smoothness.


\section{Algorithm}\label{sec:alg}

As demonstrated in our counterexamples, the Lipschitz smoothness assumption does not always hold in problem (P). In such non-Lipschitz scenarios, gradients can change drastically, posing a significant challenge for these algorithms. To address this issue, we propose a new stepsize strategy, detailed in Algorithm \ref{alg} and section \ref{section:main}, to improve performance under these challenging conditions. This strategy aims to choose the stepsize to accommodate the variance and instability in gradients, thereby enhancing the robustness of the optimization process.

In this algorithm, we start with an initial point $\tilde{w}_0$. During each iteration $t\in[T]$, either all the samples are shuffled, or we keep the order of the samples as in the last epoch. This reshuffling introduces variance in the order of samples, which can help mitigate issues related to gradient instability. For each step $j\in[n]$, we use the gradient from a single sample with number $\pi^{(t)}_j$ to update the weights $w$. The notation $\pi^{(t)}_j$ is used to denote the $j$-th element of the permutation $\pi^{(t)}$ for $j \in [n]$. Each outer loop through the data is counted as an epoch, and our convergence analysis focuses on the performance after the completion of each full epoch. 

There are multiple strategies to determine $\pi^{(t)}$:
\begin{itemize}
    \item If $\pi^{(t)}$ is a fixed permutation of $[n]$, Algorithm \ref{alg:shuffle} functions as an incremental gradient method. This method maintains a consistent order of samples, which can simplify the analysis and implementation.
    \item If $\pi^{(t)}$ is shuffled only once in the first iteration and then used in every subsequent iteration, Algorithm 1 operates as a shuffle-once algorithm. This strategy introduces randomness at the beginning but maintains a fixed order thereafter, providing a balance between randomness and stability.
    \item If $\pi^{(t)}$ is regenerated in every single iteration, Algorithm \ref{alg:shuffle} becomes a random reshuffling algorithm. This approach maximizes the randomness in the sample order, potentially offering the most robustness against the erratic behavior of non-Lipschitz gradients by constantly changing the sample order.
\end{itemize}

\begin{algorithm}[t]
\caption{Shuffling-type Gradient Algorithm}\label{alg:shuffle}
\label{alg}
\begin{algorithmic} % The number tells where the line numbering should start
\STATE \textbf{Initialization}: Choose an initial point $\tilde{w}_0\in$ $dom (F)$.
\FOR{$t=1,2,\cdots,T$}
    \STATE Set $w^{(t)}_0:=\tilde{w}_{t-1}$;
    \STATE Generate permutation $\pi^{(t)}$ of $[n]$.
    \STATE Compute non-increasing stepsize $\eta_t$.
    \FOR{$j=1,\cdots,n$}
        \STATE Update $w^{(t)}_j:=w^{(t)}_{j-1}-\frac{\eta_t}{n}\nabla f(w^{(t)}_{j-1};\pi^{(t)}_j).$
    \ENDFOR
    \STATE Set $\tilde{w}_t:=w^{(t)}_n.$
\ENDFOR 
\end{algorithmic}
\end{algorithm}

Although the random reshuffling scheme is most used in practice, each of these strategies offers distinct advantages and can be selected based on the specific requirements and characteristics of the optimization problem at hand. For this reason, we will give convergence rates for random and arbitrary shuffling scheme.

\section{Convergence Analysis}
\label{section:main}

\subsection{Main Results}

In this section, we present the main results of our convergence analysis. Our findings indicate that, with proper stepsizes, it is possible to achieve the same convergence rate, up to a logarithm difference, as under the Lipschitz smoothness assumption. First, we introduce the following assumptions regarding problem (P). Assumption \ref{assumption:first} is a standard assumption, and Assumption \ref{assumption:lsmooth} requires all $F$ and $f(\cdot;i)$ to be $\ell$-smooth.

\begin{assumption}
\label{assumption:first}
    $\mathrm{dom} (F):=\{w\in\mathbb{R}^d:F(w)<+\infty\}\neq \emptyset$ and $F^*:=\inf_{w\in\mathbb{R}^d} F(w)> -\infty.$
\end{assumption}

\begin{assumption}
\label{assumption:lsmooth}
$F$ and $f(\cdot;i)$ are $\ell$-smooth for some sub-quadratic function $\ell$, $\forall i\in[n]$.
\end{assumption}

Here, we assume all functions share the same $\ell$ function without loss of generality, as we can always choose the pointwise maximum of all their $\ell$ functions. We define $p$ to be the degree of the $\ell$ function such that $p = \text{sup}_{p\geq 0} \{p|\lim_{w\rightarrow\infty}\frac{\ell(w)}{w^p}>0\}.$
Since $\ell$ is sub-quadratic, we have $0\leq p < 2$. 

Next, we introduce our assumption about the gradient variances.

% \begin{assumption}
% \label{assumption:variance}
% There exist a constant $\sigma\in(0,+\infty)$ such that $\forall i\in [n]$,
% \begin{equation}
% \left\| \nabla f(w; i) - \nabla F(w) \right\| \leq \sigma, \; a.s., \; \forall w \in \text{dom} (F).
% \end{equation}
% \end{assumption}

% \begin{assumption}
%     \label{assumption:subgaussian}
%     There exists a constant $R>0$ such that if $i\in[n]$ is uniformly distributed,
%     \begin{equation}
%         \mathbb{P}(\|\nabla f(w;i)-\nabla F(w)\|>s)\leq 2\exp(-\frac{s^2}{2R^2}), \; \forall s>0, \;\forall w\in \text{dom} (F).
%     \end{equation}
% \end{assumption}

\begin{assumption}
    \label{assumption:varianceexpectation}
There exist two constants $\sigma, A\in(0,+\infty)$ such that $\forall i\in [n]$,
\begin{equation}
\frac{1}{n}\sum_{i=1}^n\left\| \nabla f(w; i) - \nabla F(w) \right\|^2 \leq A\|\nabla F(w)\|^2+\sigma^2, 
\end{equation}
$a.s., \; \forall w \in \text{dom} (F).$
\end{assumption}

Since component gradients behave as gradient estimations of the full gradient, this assumption can be viewed as a generalization of the more common assumption $\mathbb{E}[|\nabla F(w;\xi)-\nabla F(w)|]\leq \sigma^2$, which is used in most $\ell$-smooth work.

% We include different variance assumptions to show the change of complexity under different circumstances. Assumptions \ref{assumption:variance} and \ref{assumption:subgaussian} are stronger than the most standard bounded variance assumption $\mathbb{E}[\|\nabla f(w;i)-\nabla F(w)\|^2]\leq \sigma^2$ for some $\sigma >0$. Among them, the first one is stronger, requiring the variance to be almost bounded, while the second one allows for an exponential proportion of exceptions, making Gaussian distribution possible. These are common assumptions used in papers addressing non-Lipschitz smoothness, such as \citep{zhang2019gradient}, \citep{zhang2020improved}, and \citep{li2023convergence}. Extending to the bounded variance assumption \ref{assumption:varianceexpectation} does not change the dependency of complexity on $\epsilon$, but it does increase the dependency on $n$, which will be elaborated in Corollary \ref{theorem:originalassumption}.

\subsubsection{Nonconvex Case}
\label{section:nonconvex}

Let us denote $\Delta_1:=F(w^{(1)}_0)-F^*$. Under Assumptions \ref{assumption:first} to \ref{assumption:varianceexpectation}, we have the following result for random shuffling scheme. Proofs can be found in Appendix \ref{appendix:lemmas} and \ref{appendix:proofsfornonconvex}.

% \begin{theorem}
% \label{theorem:nonconvexconstant}
%     Suppose Assumptions \ref{assumption:first}, \ref{assumption:lsmooth} and \ref{assumption:variance} hold. Let $\{\tilde{w_t}\}_{t=1}^T$ be generated by Algorithm 1 with random reshuffling scheme. For any $0<\delta<1$, we denote $H:=\frac{4\Delta_1}{\delta}$, $G:=\sup\{u\geq 0|u^2\leq 2\ell(2u)\cdot H\}<\infty$, $G':=G+\sigma$, $L:=\ell(2G')$. Choose $\eta_t$ and $T$ such that
%     $$\eta_t\leq \frac{1}{2L}, \; \sum_{t=1}^T \eta_t^3 \leq \frac{n\Delta_1}{L^2\sigma^2}, \; T\geq \frac{32\Delta_1}{\eta_T \delta\epsilon^2},$$ then with probability at least $1-\delta$, we have $\|\nabla F(w^{(t)}_0)\|\leq G$ for every $1\leq t\leq T$ and $$\frac{1}{T}\sum_{t=1}^T\|\nabla F(w^{(t)}_0)\|^2\leq \epsilon^2.$$
% \end{theorem}

% By choosing $\eta_t=\eta = \mathcal{O}\left(\sqrt[3]{\frac{n}{T}}\right)$, we achieve a complexity of $\mathcal{O}\left(\frac{1}{\sqrt{n} \epsilon^3}\right)$ outer iterations and $\mathcal{O}\left(\frac{\sqrt{n}}{\epsilon^3}\right)$ total number of gradient evaluations, ignoring constants. This result aligns with the convergence rate for the random reshuffling scheme in \citep{nguyen2021unified}, with relaxed smoothness assumptions. We will further compare the complexity under Corollary \ref{theorem:originalassumption}.

% If we assume \ref{assumption:subgaussian} instead of \ref{assumption:variance}, we have the following result.

% \begin{corollary}
% \label{theorem:subgaussiannonconvexrr}
%     Suppose Assumptions \ref{assumption:first}, \ref{assumption:lsmooth} and \ref{assumption:subgaussian} hold. Let $\{\tilde{w_t}\}_{t=1}^T$ be generated by Algorithm 1 with random reshuffling scheme. For any $0<\delta<1$, we denote $\sigma:= R\sqrt{2\log(\frac{8nT}{\delta})}$, $H:=\frac{8\Delta_1}{\delta}$, $G:=\sup\{u\geq 0|u^2\leq 2\ell(2u)\cdot H\}<\infty$, $G':=G+\sigma$, $L:=\ell(2G')$. Choose $\eta_t$ and $T$ such that
%     $$\eta_t\leq \frac{1}{2L}, \; \sum_{t=1}^T \eta_t^3 \leq \frac{n\Delta_1}{L^2\sigma^2}, \; T\geq \frac{32\Delta_1}{\eta_T \delta\epsilon^2},$$ then with probability at least $1-\delta$, we have $\|\nabla F(w^{(t)}_0)\|\leq G$ for every $1\leq t\leq T$ and $$\frac{1}{T}\sum_{t=1}^T\|\nabla F(w^{(t)}_0)\|^2\leq \epsilon^2.$$
% \end{corollary}

% In this version of the theorem, we change the variance assumption to a sub-Gaussian distribution. Consequently, $\sigma$ here is $\mathcal{O}(\sqrt{\log(nT)})$ instead of a constant, which influences the total complexity. In this case, we need $\mathcal{O}\left(\frac{1}{\sqrt{n}\epsilon^3}\log^{\frac{p+1}{2}}\left(\frac{n}{\epsilon}\right)\right)$ outer iterations, where $p$ is the degree of the $\ell$ function in Assumption \ref{assumption:lsmooth}. Since $\ell$ is sub-quadratic, we have $\frac{p+1}{2} < \frac{3}{2}$. The complexity is slightly worse than the almost surely bounded variance case by a polynomial factor of the logarithm function of $n/\epsilon$.

% In the corollary below we state the result if we relax the assumption to bounded variance expectation. The proof can be found in Appendix \ref{appendix:proofsfornonconvex}.

\begin{theorem}
\label{theorem:originalassumption}
    Suppose Assumptions \ref{assumption:first}, \ref{assumption:lsmooth} and \ref{assumption:varianceexpectation} hold, Let $\{\tilde{w_t}\}_{t=1}^T$ be generated by Algorithm \ref{alg:shuffle} with random reshuffling scheme. For any $0<\delta<1$, we denote $H:=\frac{4\Delta_1}{\delta}$, $G:=\sup\{u\geq 0|u^2\leq 2\ell(2u)\cdot H\}$, $G':=\sqrt{2(1+nA)}G+\sqrt{2n}\sigma$, $L:=\ell(2G')$. For any $\epsilon>0$, choose $\eta_t$ and $T$ such that
    $$\eta_t\leq \frac{1}{2L\sqrt{\frac{A}{n}+1}}, \; \sum_{t=1}^T \eta_t^3 \leq \frac{n\Delta_1}{L^2\sigma^2}, \; T\geq \frac{32\Delta_1}{\eta_T \delta\epsilon^2},$$ then with probability at least $1-\delta$, we have $\|\nabla F(w^{(t)}_0)\|\leq G$ for every $1\leq t\leq T$
    $$\frac{1}{T}\sum_{t=1}^T\|\nabla F(w^{(t)}_0)\|^2\leq \epsilon^2.$$
\end{theorem}

\begin{remark}
    By choosing $\eta_t=\eta=\mathcal{O}(\sqrt[3]{\frac{n^{1-p}}{T}})=\mathcal{O}(n^\frac{1-p}{2}\epsilon)$, we can achieve a complexity of $T=\mathcal{O}(\frac{n^\frac{p-1}{2}}{\epsilon^3})$ outer iterations and $\mathcal{O}(\frac{n^\frac{p+1}{2}}{\epsilon^3})$ total number of gradient evaluations, ignoring constants, where $p$ is the order of the $\ell$ function in Definition \ref{def:gs}. 
    As $p$ goes to $0$, $\ell$-smoothness degenerates to the traditional Lipschitz smoothness, and our total number of gradient evaluations goes to $\mathcal{O}(\frac{\sqrt{n}}{\epsilon^3})$ once again, which matches the complexity in Corollary 1 of \citet{nguyen2021unified}. If $\epsilon\leq 1/\sqrt{n}$, one possible stepsize is $\eta=\frac{\sqrt{n}\epsilon}{2L\sqrt{\frac{A}{n}+1}}.$
\end{remark}

Our result here has polynomial dependency on $\frac{1}{\delta}$, $T=\mathcal{O}(\delta^{-\frac{3}{2}-\frac{p}{2-p}})$.
% , which contrasts with many results for shuffling gradient algorithms under Lipschitz smoothness assumptions, e.g. theorem 1 in \citet{ahn2020sgd}. 
It is important to note that, in our setting, $\delta$ accounts for the probability that Lipschitz smoothness does not hold—a consideration absent in standard Lipschitz smoothness settings. Therefore, the dependency here is not as good as in $L$-smoothness cases. In fact, a polynomial dependency on $\delta$ is typical in papers with similar smoothness assumptions, e.g. in \citet{li2023convex}, \citet{li2023convergence}, \citet{xiandelving} and \citet{zhang2024convergence}. 

Next we consider arbitrary $\pi^{(t)}$ scheme in Algorithm \ref{alg:shuffle}.

\begin{theorem}
\label{theorem:nonconvexallscheme}
    Suppose Assumptions \ref{assumption:first}, \ref{assumption:lsmooth} and \ref{assumption:varianceexpectation} hold. Let $\{\tilde{w_t}\}_{t=1}^T$ be generated by Algorithm \ref{alg:shuffle} with arbitrary scheme. Define $H=2\Delta_1$, $G:=\sup\{u\geq 0|u^2\leq 2\ell(2u)\cdot H\}$, $G':=\sqrt{2(1+nA)}G+\sqrt{2n}\sigma$, $L:=\ell(2G')$. For any $\epsilon>0$, choose $\eta_t$ and $T$ such that
    $$\eta_t\leq \frac{1}{L\sqrt{2(3A+2)}}, \; \sum_{t=1}^T \eta_t^3 \leq \frac{2\Delta_1}{3\sigma^2L^2}, \; T\geq \frac{8\Delta_1}{\eta_T \epsilon^2},$$ then we have $\|\nabla F(w^{(t)}_0)\|\leq G$ for every $1\leq t\leq T$ and $$\frac{1}{T}\sum_{t=1}^T\|\nabla F(w^{(t)}_0)\|^2\leq \epsilon^2.$$
\end{theorem}

This theorem gives the convergence rate for arbitrary scheme in Algorithm \ref{alg:shuffle}. By choosing $\eta_t = \eta = \mathcal{O}\left(\sqrt[3]{\frac{1}{n^p T}}\right)=\mathcal{O}\left(\frac{\epsilon}{n^\frac{p}{2}}\right)$, we achieve a complexity of $\mathcal{O}\left(\frac{n^{\frac{p}{2}}}{\epsilon^3}\right)$ outer iterations and $\mathcal{O}\left(\frac{n^{\frac{p}{2}+1}}{\epsilon^3}\right)$ total gradient evaluations, ignoring constants. Without the randomness in $\pi$ in every iteration, the complexity's dependency on $n$ is increased by $\mathcal{O}(\sqrt{n})$. One possible stepsize is $\eta=\frac{\epsilon}{L\sqrt{2(3A+2)}}.$

% \begin{remark}
%     If we replace Assumption \ref{assumption:variance} with the sub-Gaussian Assumption \ref{assumption:subgaussian}, the complexity will increase by a polynomial logarithm function of $n/\epsilon$. This result also applies to the convergence results in the convex cases but we omit them for simplicity.
% \end{remark}

\subsubsection{Strongly Convex Case}
\label{section:strongly}

For strongly convex case, we give results for both random reshuffling scheme and arbitrary scheme, with constant learning rate. Proof can be found in Appendix \ref{stronglyproof}.

\begin{assumption}
\label{assumption:strongcon}
    Function $F$ in (P) is $\mu$-strongly convex on $\mathrm{dom}(F)$.
\end{assumption}

\begin{theorem}
\label{theorem:stronglyconvexrr}
    Suppose Assumptions \ref{assumption:first}, \ref{assumption:lsmooth}, \ref{assumption:varianceexpectation} and \ref{assumption:strongcon} hold. Let $\{\tilde{w_t}\}_{t=1}^T$ be generated by Algorithm 1 with random reshuffling scheme. For any $0<\delta<1$,  we denote $H:=\max\{\frac{3\sigma^2}{4\mu}\log\frac{4}{\epsilon}+\Delta_1, \frac{4\Delta_1}{\delta}\}$, $G:=\sup\{u\geq 0|u^2\leq 2\ell(2u)\cdot H\}$, $G':=\sqrt{2(1+nA)}G+\sqrt{2n}\sigma$, $L:=\ell(2G')$. For any $\epsilon>0$, if we choose $\eta_t$ and $T$ such that
    $$\eta_t=\eta=\frac{4\log(\sqrt{n}T)}{\mu T}, T\geq 4\sqrt{\frac{\Delta_1}{n\delta\epsilon}},\frac{T}{\log(\sqrt{n}T)}\geq$$ 
    $$ \frac{4}{\mu}\max\left\{2,L\sqrt{2(3A+2)},L\sigma\sqrt{\frac{8}{n\mu\delta\epsilon}},\sqrt[3]{\frac{T\sigma^2 L^2}{n\Delta_1}}\right\},$$
    then for any $0<\delta<1$, with probability at least $1-\delta$ we have
    $$F(w^{(T+1)}_0)-F^*\leq \epsilon.$$
\end{theorem}

% Since $H=\Theta(\log\frac{1}{\epsilon})$, we have $G=\mathcal{O}(\log^{\frac{1}{2-p}}\frac{1}{\epsilon})$ and $L=\mathcal{O}(\log^{\frac{p}{2-p}}\frac{1}{\epsilon})$, where $0\leq p<2$ is the degree of the $\ell$ function. 
In Theorem \ref{theorem:stronglyconvexrr}, we can achieve a complexity of $\widetilde{\mathcal{O}}\left(n^{\frac{p-1}{2}}\epsilon^{-\frac{1}{2}}\right)$ outer iterations and $\widetilde{\mathcal{O}}\left(n^{\frac{p+1}{2}}\epsilon^{-\frac{1}{2}}\right)$ total gradient evaluations with $\eta=\widetilde{\mathcal{O}}\left(n^{\frac{1-p}{2}}\epsilon^{\frac{1}{2}}\right)$, ignoring constants. This matches the result in \citet{nguyen2021unified} with the same assumptions in the degenerate case of $p=0$. The dependence on $\delta$ is $T=\mathcal{O}(\delta^{-\frac{1}{2}-\frac{p}{2-p}}).$

% The inequality constraints on $T$ require $\frac{T}{\log(\sqrt{n}T)} \geq \mathcal{O}\left(\frac{L}{\mu}\right)$, where $\frac{L}{\mu}$ is the condition number of the functions in (P). Previous literature, such as \citep{mishchenko2020random} and \citep{safran2021random}, has shown similar requirements on $T$.

It is not hard to follow proof of Theorem \ref{theorem:nonconvexallscheme} for arbitrary scheme in strongly convex case and achieve a complexity of $\widetilde{\mathcal{O}}\left(n^{\frac{p}{2}+1}\epsilon^{-\frac{1}{2}}\right)$ total gradient evaluations. Here we give a slightly stronger result where we remove Assumption \ref{assumption:varianceexpectation} to match the corresponding result in Lipschitz smooth case.

\begin{theorem}
\label{theorem:stronglyconvexany}
 Suppose Assumptions \ref{assumption:first}, \ref{assumption:lsmooth} and \ref{assumption:strongcon} hold. Let $\{\tilde{w_t}\}_{t=1}^T$ be generated by Algorithm 1 with arbitrary scheme. We denote $S=\{w|F(w)\leq F(w^{(1)}_0)\}$, $G'=\max_w\{\|\nabla f(w;i)\||w\in S, i\in[n]\}$, $L:=\ell(2G')$. For any $\epsilon>0$, choose $\eta_t$ and $T$ such that
 $$\eta_t=\eta=\frac{6\log(T)}{\mu T}\leq \frac{\Delta_1\mu^2}{9(\mu^2+L^2)\sigma^2_*},$$ 
 $$T=\widetilde{\mathcal{O}}(\epsilon^{-\frac12})\geq \frac{12L^2\log(T)}{\mu^2},$$
    where $\sigma_*$ is the standard deviation at $w_*$. Then we have $\|\nabla F(w^{(t)}_0)\|\leq G'$ and 
    $$F(w^{(T+1)}_0)-F^*\leq \epsilon.$$
\end{theorem}

% \begin{theorem}
% \label{theorem:stronglyconvexany}
%  Suppose Assumptions \ref{assumption:first}, \ref{assumption:lsmooth}, \ref{assumption:varianceexpectation} and \ref{assumption:strongcon} hold. Let $\{\tilde{w_t}\}_{t=1}^T$ be generated by Algorithm 1 with arbitrary scheme. We denote $H:=\max\{\frac{3\sigma^2}{4\mu}\log\frac{4}{\epsilon}+\Delta_1, 2\Delta_1\}$, $G:=\sup\{u\geq 0|u^2\leq 2\ell(2u)\cdot H\}$, $G':=\sqrt{2(1+nA)}G+\sqrt{2n}\sigma$, $L:=\ell(2G')$. For any $\epsilon>0$, choose $\eta_t$ and $T$ such that
%  $$\eta_t=\eta=\frac{4\log(\sqrt{n}T)}{\mu T}, T\geq \sqrt{\frac{\Delta_1}{n\epsilon}},$$ 
% $$\frac{T}{\log(\sqrt{n}T)}\geq \frac{4}{\mu}\max\left\{2,L\sqrt{2(3A+2)},\frac{L\sigma}{\mu}\sqrt{\frac{3}{\mu\epsilon}},\sqrt[3]{\frac{T\sigma^2 L^2}{n\Delta_1}}\right\},$$
%  % $$\eta_t=\eta=\frac{\log T}{\mu T}, \frac{T}{\log T}\geq \frac{1}{\mu}\max\left\{\frac{2G'L}{G},\sqrt[3]{\frac{TL^2(2G^2+3\sigma^2)}{2\Delta_1}}\right\},$$
% then it holds that
% $$F(w^{(T+1)}_0)-F^*\leq \epsilon.$$
%   % \begin{equation*}
%   %   \begin{split}
%   %     F(w^{(T+1)}_0)-F^*
%   % \le  \left(F(w^{(1)}_0)-F^*\right)\frac{1}{T^2} + \frac{L^2(2G^2+3\sigma^2)}{\mu^3}\left(\frac{\log T}{T}\right)^2 = \mathcal{O}\left(\frac{\log^2 T}{T^2}\right).
%   %   \end{split}
%   % \end{equation*}
% \end{theorem}

In Theorem \ref{theorem:stronglyconvexany}, we achieve a complexity of $\widetilde{\mathcal{O}}\left(\epsilon^{-1/2}\right)$ outer iterations and $\widetilde{\mathcal{O}}\left(n\epsilon^{-1/2}\right)$ total gradient evaluations with $\eta=\widetilde{\mathcal{O}}\left(n^{-1}\epsilon^{\frac{1}{2}}\right)$, ignoring constants. It should be noted that the constant $G'$ here is implicitly determined by constant $p$ and can potentially be large. Therefore, the complexity here cannot be directly compared with those in Theorem \ref{theorem:nonconvexallscheme} or \ref{theorem:stronglyconvexrr}.
% Therefore, the randomness in $\pi^{(t)}$ of the random reshuffling scheme reduces the total complexity by an $\widetilde{\mathcal{O}}(\sqrt{n})$ term.

% \begin{theorem}
% \label{theorem:stronglyconvexany}
%  Suppose Assumptions \ref{assumption:first}, \ref{assumption:lsmooth}, \ref{assumption:varianceexpectation} and \ref{assumption:strongcon} hold. Let $\{\tilde{w_t}\}_{t=1}^T$ be generated by Algorithm 1 with arbitrary scheme. Define $H$, $G$, $G'$, $L$ and $r$ as in Theorem \ref{theorem:originalassumption}. Choose $\eta_t$ and $T$ such that
%  $$\eta_t=\eta=\frac{\log T}{\mu T},$$ 
%  $$\frac{T}{\log T}\geq \frac{1}{\mu}\max\left\{\frac{2G'L}{G},\sqrt[3]{\frac{TL^2(2G^2+3\sigma^2)}{2\Delta_1}}\right\},$$
% then it holds that
%   \begin{equation*}
%     \begin{split}
%       F(w^{(T+1)}_0)-F^*
%   \le  \left(F(w^{(1)}_0)-F^*\right)\frac{1}{T^2} + \frac{L^2(2G^2+3\sigma^2)}{\mu^3}\left(\frac{\log T}{T}\right)^2 = \mathcal{O}\left(\frac{\log^2 T}{T^2}\right).
%     \end{split}
%   \end{equation*}
% \end{theorem}

% In this setting, $H$ is a constant. Therefore, in Theorem \ref{theorem:stronglyconvexany}, we achieve a complexity of $\widetilde{\mathcal{O}}\left(\epsilon^{-1/2}\right)$ outer iterations and $\widetilde{\mathcal{O}}\left(n\epsilon^{-1/2}\right)$ total gradient evaluations, ignoring constants. Therefore, the randomness in $\pi^{(t)}$ of the random reshuffling scheme reduces the total complexity by an $\widetilde{\mathcal{O}}(\sqrt{n})$ term.

\subsubsection{Non-strongly Convex Case}
\label{section:nonstrongly}

Next we consider the case where only non-strongly convexity are assumed. In the following theorem, we assume the optimal solution exists and denote one as $w_*$, the standard deviation at $w_*$ as $\sigma_*:=\sqrt{\frac{1}{n}\sum_{i=1}^n\|\nabla f(w_*;i)\|^2}$ and the average value of $\{w^{(t)}_0\}_{t=1}^T$ as $\bar{w}_{T}=\frac{1}{T}\sum_{t=1}^{T}w^{(t)}_0$. Proof for this section can be found in Appendix \ref{convexproof}.

\begin{assumption}
\label{assumption:nonstrongly}
    Functions $f(\cdot;i)$ in (P) are convex on$\mathrm{dom}(F)$, for all $i\in[n]$.
\end{assumption}

\begin{theorem}
\label{theorem:nonstrongly}
Suppose Assumptions \ref{assumption:first}, \ref{assumption:lsmooth}, \ref{assumption:varianceexpectation} and \ref{assumption:nonstrongly} hold. Let $\{\tilde{w_t}\}_{t=1}^T$ be generated by Algorithm 1 with random reshuffling scheme. For any $0<\delta<1$, define $H$, $G$, $G'$, $L$ as in Theorem \ref{theorem:originalassumption}. For any $\epsilon>0$, choose $\eta_t=\eta$ and $T$ such that
    $$\eta\leq \min \left\{\frac{1}{2L\sqrt{\frac{A}{n}+1}}, \sqrt[3]{\frac{n\Delta_1}{T\sigma^2 L^2}}, \sqrt[3]{\frac{3n\|w^{(1)}_0-w_*\|^2}{2LT\sigma_*^2}}\right\},$$
    $$T\geq \frac{4\|w^{(1)}_0-w_*\|^2}{\eta \delta\epsilon},$$
    then with probability at least $1-\delta$, we have $\|\nabla F(w^{(t)}_0)\|\leq G$ for every $1\leq t\leq T$ and $$F(\bar{w}_T)-F^*\leq \epsilon.$$
\end{theorem}

By choosing $\eta = \mathcal{O}\left(\sqrt[3]{\frac{n^{1-p}}{T}}\right)=\mathcal{O}\left(n^\frac{1-p}{2}\epsilon^{0.5}\right)$, we achieve a complexity of $\mathcal{O}\left(\frac{n^{\frac{p-1}{2}}}{\epsilon^{1.5}}\right)$ outer iterations and $\mathcal{O}\left(\frac{n^{\frac{p+1}{2}}}{\epsilon^{1.5}}\right)$ total number of gradient evaluations, ignoring constants. The dependency on $\delta$ is $T=\mathcal{O}(\delta^{-\frac{3}{2}-\frac{p}{2-p}})$. If $\epsilon\leq 1/n$, one possible stepsize is $\eta=\frac{\sqrt{n\epsilon}}{2L\sqrt{\frac{A}{n}+1}}.$

Similarly, we can follow proof of Theorem \ref{theorem:nonconvexallscheme} for arbitrary scheme and achieve a complexity of $\mathcal{O}\left(\frac{n^{\frac{p}{2}+1}}{\epsilon^{1.5}}\right)$ total gradient evaluations. Now we give a result without variance assumption \ref{assumption:varianceexpectation}.

\begin{theorem}
\label{theorem:nonstronglyany}
    Suppose Assumptions \ref{assumption:first}, \ref{assumption:lsmooth} and \ref{assumption:nonstrongly} hold. Let $\{\tilde{w_t}\}_{t=1}^T$ be generated by Algorithm 1 arbitrary scheme. Define $S=\{w|F(w)\leq F(w^{(1)}_0)\}$, $G'=\max_w\{\|\nabla f(w;i)\||w\in S, i\in[n]\}<\infty$, $L=\ell(2G')$. For any $\epsilon>0$, choose $\eta_t=\eta$ and $T$ such that
    $$\eta\leq \frac{1}{G'}\sqrt\frac{3\epsilon}{2L}, T=\mathcal{O}(\epsilon^{-1.5})\geq \frac{\|w^{(1)}_0-w_*\|^2}{\eta \epsilon},$$
    then we have $\|\nabla F(w^{(t)}_0)\|\leq G$ for every $1\leq t\leq T$ and
    $$\min_{t\in[T]}F(w_T)-F^*\leq \epsilon.$$
\end{theorem}

% \begin{theorem}
% \label{theorem:nonstronglyany}
%     Suppose Assumptions \ref{assumption:first}, \ref{assumption:lsmooth}, \ref{assumption:varianceexpectation} and \ref{assumption:nonstrongly} hold. Let $\{\tilde{w_t}\}_{t=1}^T$ be generated by Algorithm 1 arbitrary scheme. Define $H$, $G$, $G'$, $L$ as in Theorem \ref{theorem:nonconvexallscheme}. For any $\epsilon>0$, choose $\eta_t$ and $T$ such that
%     $$\eta_t=\eta\leq \min \left\{\frac{1}{L\sqrt{2(3A+2)}}, \sqrt[3]{\frac{2\Delta_1}{TL^2(2G^2+3\sigma^2)}},\sqrt[3]{\frac{\|w^{(1)}_0-w_*\|^2}{LT\sigma^2}} \right\},$$
%     $$T\geq \frac{\|w^{(1)}_0-w_*\|^2}{\eta \epsilon},$$
%     then we have $\|\nabla F(w^{(t)}_0)\|\leq G$ for every $1\leq t\leq T$ and
%     $$F(\bar{w}_T)-F^*\leq \epsilon.$$
% \end{theorem}

By choosing $\eta=\mathcal{O}(\sqrt[3]{\frac{1}{T}})$, we have the complexity of $\mathcal{O}(\frac{1}{\epsilon^{1.5}})$ outer iterations and $\mathcal{O}(\frac{n}{\epsilon^{1.5}})$ total number of gradient evaluations, ignoring constants. One possible stepsize is $\eta=\frac{\sqrt{3\epsilon}}{G'\sqrt{L}}.$

% \subsubsection{Convex Case without Variance Assumption}

% In this section, we analyze the performance of Algorithm \ref{alg:shuffle} without any assumptions about variance in strongly and non-strongly convex cases. Proof can be found in Appendix \ref{appendix:final}.

% \begin{theorem}
% \label{theorem:noassum1}
%     Suppose Assumptions \ref{assumption:first}, \ref{assumption:lsmooth} and \ref{assumption:strongcon} hold. Define $S=\{w|F(w)\leq F(w^{(1)}_0)\}$, $G'=\max_w\{\|\nabla f(w;i)\||w\in S, i\in[n]\}$, $L=\ell (2G')$. If we choose $\eta_t$ and $T$ such that 
%     $$\eta_t=\eta=\frac{6\log(T)}{\mu T}\leq \frac{\Delta_1\mu^2}{9(\mu^2+L^2)\sigma^2_*}, T=\widetilde{\mathcal{O}}(\epsilon^{-\frac12})\geq \frac{12L^2\log(T)}{\mu^2},$$
%     where $\sigma_*$ is the standard deviation at $w_*$. Then we have $\|\nabla F(w^{(t)}_0)\|\leq G'$ and 
%     $$F(w^{(T+1)}_0)-F^*\leq \epsilon.$$
% \end{theorem}

% \begin{theorem}
% \label{theorem:noassum2}
%     Suppose Assumptions \ref{assumption:first}, \ref{assumption:lsmooth} and \ref{assumption:nonstrongly} hold.
% \end{theorem}

\subsection{Proof Sketch and Technical Novelty}

Broadly speaking, our approach involves two main goals: first, demonstrating that Lipschitz smoothness is maintained with high probability along the training trajectory $\{\tilde{w}_t\}$, and second, showing that, conditioned on Lipschitz smoothness, the summation of gradient norms is bounded with high probability. Here we slightly abuse the term 'Lipschitz' and 'Lipschitz smoothness' to refer to the property between neighboring steps along the training trajectory.

For the first goal, in Lemma \ref{Lemma:epsilon}, we prove by induction that when starting an iteration with a bounded gradient, the entire training trajectory during this iteration will have bounded gradients. Consequently, we only need to verify the Lipschitz smoothness condition at the start of each iteration. However, at this point, the two goals become intertwined. We need Lipschitz smoothness to bound the gradient differences, but we also need the gradient norm bounds to establish Lipschitz smoothness. Our solution is to address both issues simultaneously.

Assuming that, before a stopping time $\tau$, Lipschitz smoothness holds, we bound the gradient norm up to that time in Lemma \ref{Lemma:Ftau}. However, this process is nontrivial. Since we are examining behavior before a stopping time, every expectation is now conditioned on $t < \tau$, rendering all previous estimations for shuffling gradient algorithms inapplicable. This presents a contradiction: we want to condition on $t < \tau$ when applying Lipschitz smoothness, but we do not want this condition when estimating other quantities. In Lemma \ref{Lemma:Ftau}, we find a method to separately handle these two requirements, allowing us to achieve both goals simultaneously.


% For the first part, we aim to show that the function values along the trajectory $\{F(\tilde{w}_t)\}$ remain bounded with high probability. This follows from the definition of $\ell$-smoothness and Lemma \ref{Lemma:gradienttodiff}. To estimate the sequence of function values, it is necessary to bound the difference between the true gradient and the gradient computed from a single sample. During this estimation process, two conditions are frequently needed: Lipschitz smoothness along the trajectory and the uniform distribution of $\pi^{(t)}$. However, these two conditions can be in conflict: assuming Lipschitz smoothness means that $\pi^{(t)}$ is no longer uniformly distributed. Our proof strategy addresses these two conditions separately, thereby avoiding the conflict and ensuring a consistent approach in.

% For the second part, once we establish Lipschitz smoothness along the trajectory with high probability, we can estimate the behavior of the gradient norm or the optimality gap based on the algorithm's performance up to the first instance when Lipschitz smoothness is violated. This involves estimating the summation of the desired quantity up to the point of the first violation of Lipschitz smoothness.

\subsection{Limitations and Future works}

Although we have proved upper bounds for the complexity of shuffling gradient algorithms, there are certain limitations in our work that we leave for future research:

\begin{itemize}
    \item First, as is common with many optimization algorithms, it is challenging to verify that the bounds presented are indeed the lower bounds. Future work could explore improving these results, for instance, by reducing the dependency on $\delta$ to a logarithmic factor, or by proving that the current bounds are, in fact, tight lower bounds.
    \item Second, although we showed results for arbitrary shuffling schemes, there are better results for single shuffling under Lipschitz smoothness, for example \citet{ahn2020sgd} proved $\mathcal{O}(\frac{1}{nT^2})$ convergence rate for strongly convex objectives. It is interesting to see whether we can achieve the same convergence rate with $\ell$-smoothness as well.
    \item The results in Theorem \ref{theorem:stronglyconvexany} and Theorem \ref{theorem:nonstronglyany} depends on constant $G'$ that can be potentially very large and hard to verify. In the absence of both Lipschitz smoothness and bounded variance, the behavior of gradients can be hard to track. We hope our results here can be a first step for future work.
    \item Lastly, shuffling gradient methods have been integrated with variance reduction techniques \citep{malinovsky2023random}. Exploring the performance of these algorithms under relaxed smoothness assumptions is another promising direction for future work.
\end{itemize}

\section{Numerical Experiments}
We compare reshuffling gradient algorithm (Algorithm \ref{alg:shuffle}) with SGD on multiple $\ell$-smooth optimization problems to prove its effectiveness. Experiments are conducted with different shuffling schemes, on convex, strongly convex and nonconvex objective functions, including synthetic functions, phase retrieval, distributionally robust optimization (DRO) and image classification. 
% The whole experiment runs on Python 3.9 in a MacBook Pro laptop with 500 GB Storage and 8-core CPU (16 GB Memory). 
\subsection{Convex and Strongly Convex Settings}
We first consider convex functions $f_{i,k}(x)=x_i^4+kx_i$  of $x\in\mathbb{R}^{50}$ for all $(i,k)\in \mathcal{E}:=\{1,2,\ldots,50\}\times\{-10,-9,\ldots,9,10\}$, as well as their sample average $f(x)=\frac{1}{1050}\sum_{(k,i)\in\mathcal{E}}f_{i,k}(x)=\frac{1}{50}\sum_{i=1}^{50} x_i^4$. It can be easily verified that $f$ and all $f_{i,k}$ are convex but not strongly convex, and $\ell$-smooth (with $\ell(u)=3u^{2/3}$) but not Lipschitz-smooth. Then we compare reshuffling gradient algorithm (Algorithm \ref{alg:shuffle}) with SGD on the objective $\min_{x\in\mathbb{R}^{50}}f(x)$. Specifically, for each SGD update $x\leftarrow x-\eta\nabla f_{k,i}(x)$, $(k,i)\in\mathcal{E}$ is obtained uniformly at random. For Algorithm \ref{alg:shuffle}, we adopt three shuffling schemes as elaborated in Section \ref{sec:alg}. The fixed-shuffling scheme and shuffling-once fix all permutations $\pi^{(t)}$ respectively to be the natural sequence $(1,-10), (1,-9),\ldots (50,10)$ and its random permutation at the beginning, while the uniform-shuffling scheme obtains permutations $\pi^{(t)}$ uniformly at random and independently for all iterations $t$. We implement each algorithm 100 times with initialization $x_0=[1,\ldots,1]$ and fine-tuned stepsizes 0.01 (i.e., $\eta=0.01$ for SGD and $\frac{\eta_t}{n}=0.01$ for Algorithm \ref{alg:shuffle}), which takes around 3 minutes in total. We plot the learning curves of $f(x_t)$ averaged among the 100 times, as well as the 95\% and 5\% percentiles in the left of Figure \ref{fig_convex}, which shows that Algorithm \ref{alg:shuffle} with all shuffling schemes converges faster than SGD. 

Then we consider strongly convex functions $f_{j,k}(x)=\exp(x_j-k)+\exp(k-x_j)+\frac{1}{2}||x||^2$ for $(i,k)\in\mathcal{E}$ and their sample average below. 
\begin{align*}
    &f(x)=\frac{1}{1050}\sum_{(k,i)\in\mathcal{E}}f_{i,k}(x)=\frac{1}{2}||x||^2+\\
    &\frac{\exp(n+1)-\exp(-n)}{1050[\exp(1)-1]}\sum_{j=1}^{50} [\exp(x_j)+\exp(-x_j)].
\end{align*}
All these functions $f_{j,k}$ and $f$ are 1-strongly convex and $\ell$-smooth (with $\ell(u)=5u+5$) but not Lipschitz-smooth. We repeat the experiment in the same procedure above, except that all the stepsizes are fine-tuned to be $10^{-5}$. The result is shown in the right of Figure \ref{fig_convex}, which also shows that Algorithm \ref{alg:shuffle} with all shuffling schemes converges faster than SGD. 

\begin{figure}[t]
\begin{minipage}{.47\textwidth}
    \centering
    \includegraphics[width=.8\textwidth]{poly.png}	
\end{minipage} 
\begin{minipage}{.47\textwidth}
    \centering
    \includegraphics[width=.8\textwidth]{exponential.png}
\end{minipage} 
\caption{Experimental Results on Convex (up) and Strongly-convex (down) Objective Functions.}\label{fig_convex}
\vspace{-1pt}
\end{figure}

\subsection{Application to Phase Retrieval and DRO}
We compare SGD with Algorithm \ref{alg:shuffle} on phase retrieval and distributionally robust optimization (DRO), which are  $\ell$-smooth but not Lipschitz smooth. We use similar setup as in \citep{chen2023generalized}. 

In the phase retrieval problem \eqref{obj_phase}, we select $m=3000$ and $d=100$, and generate independent Gaussian variables $x, a_r\sim \mathcal{N}(0,0.5I_d)$, initialization $z_0\sim\mathcal{N}(5,0.5I_d)$, as well as $y_i=|a_r^{\top}z|^2+n_i$ with noise $n_i\sim\mathcal{N}(0,4^2)$ for $i=1,...,m$. We select constant stepsizes $2\times 10^{-6}$ and $\eta_j^{(t)}\equiv \frac{0.007}{m}$ for SGD and Algorithm \ref{alg:shuffle} respectively by fine-tuning and implement each algorithm 100 times. % which takes around 17 minutes in total. 
For Algorithm \ref{alg:shuffle}, we adopt three shuffling schemes as elaborated in Section \ref{sec:alg}. The fixed-shuffling scheme and shuffling-once fix all permutations $\pi^{(t)}$ respectively to be the natural sequence $1,2,\ldots,3000$ and its random permutation at the beginning, while the uniform-shuffling scheme obtains permutations $\pi^{(t)}$ uniformly at random and independently for all iterations $t$. We plot the learning curves of the objective function values averaged among the 100 times, as well as the 95\% and 5\% percentiles in the left of Figure \ref{fig_Phase_DRO}, which shows that Algorithm \ref{alg:shuffle} with shuffle-once and uniform-shuffling schemes converge faster than SGD. 
\begin{figure}[t]
\begin{minipage}{.47\textwidth}
    \centering
    \includegraphics[width=.7\textwidth]{Phase.png}	
\end{minipage} 
\begin{minipage}{.47\textwidth}
    \centering
    \includegraphics[width=.7\textwidth]{DRO.png}
\end{minipage} 
\caption{Experimental Results on Phase Retrieval (up) and DRO (down).}\label{fig_Phase_DRO}
\vspace{-1pt}
\end{figure}

In the DRO problem \eqref{eq:DRO}, we select $\lambda=0.01$ and $\psi^*(t)=\frac{1}{4}(t+2)_+^2-1$ (corresponding to $\psi$ being $\chi^2$ divergence). For the stochastic samples $\xi$, we use the life expectancy data\footnote{\url{https://www.kaggle.com/datasets/kumarajarshi/life-expectancy-who?resource=download}} designed for regression task between the life expectancy (target) and its factors (features) of 2413 people, and preprocess the data by filling the missing values with the median of the corresponding features, censorizing and normalizing all the features \nopagebreak\footnote{The detailed process of filling missing values and censorization: \url{https://thecleverprogrammer.com/2021/01/06/life-expectancy-analysis-with-python/}}, removing two categorical features (``country'' and ``status''), and adding standard Gaussian noise to the target to get robust model. We use the first 2000 samples $\{x_i,y_i\}_{i=1}^{2000}$ with features $x_i\in \mathbb{R}^{34}$ and targets $y_i\in \mathbb{R}$ for training. We use the loss function $\ell_{\xi}(w)=\frac{1}{2}(y_{\xi}-x_{\xi}^{\top}w)^2+0.1\sum_{j=1}^{34}\ln\big(1+|w^{(j)}|)\big)$ of $w=[w^{(1)};\ldots;w^{(34)}]\in\mathbb{R}^{34}$ for any sample $x_{\xi},y_{\xi}$. We use initialization $\eta_0=0.1$ and $w_0\in\mathbb{R}^{34}$ from standard Gaussian distribution. 

Then similar to phase retrieval, we implement both SGD and the three sampling schemes of Algorithm \ref{alg:shuffle} 100 times with stepsizes $\eta_j^{(t)}\!\!=\!\frac{\eta_t}{n}\!\!=\!10^{-7}$.%, which takes 11-12 hours. 
We evaluate $\Psi(x_t)\!:=\!\min_{\eta\in\mathbb{R}} L(x_t,\eta)$ every 10 iterations. The average, 5\% and 95\% percentiles of $\Psi(x_t)$ among the 100 implementations are plotted in the right of Figure \ref{fig_Phase_DRO}, which shows that Algorithm \ref{alg:shuffle} with fixed shuffling converges faster than SGD. 

\begin{figure}
\begin{minipage}{.5\textwidth}
    \centering
    \includegraphics[width=.8\textwidth]{TrainAccuracy.png}
\end{minipage} 
\begin{minipage}{.5\textwidth}
    \centering
    \includegraphics[width=.8\textwidth]{TestAccuracy.png}
\end{minipage} 
	\caption{Experimental Results on Cifar 10 Dataset.}
	\label{fig:Cifar10}
	\vspace{-2pt}
\end{figure}


\subsection{Application to Image Classification}
We train Resnet18 \citep{he2016deep} with cross-entropy loss for image classification task on Cifar 10 dataset \citep{krizhevsky2009learning}, using SGD and Algorithm \ref{alg:shuffle} with three shuffling schemes. We implement each algorithm 100 times with batchsize 200 and stepsize $10^{-3}$. After every 250 iterations, we evaluate the sample-average loss value as well as classification accuracy on the whole training dataset and test dataset. The average, 5\% and 95\% percentiles of these evaluated metrics among the 100 implementations are plotted in Figure \ref{fig:Cifar10}, which shows that Algorithm \ref{alg:shuffle} with fixed-shuffling scheme outperforms SGD on both training and test data, and Algorithm \ref{alg:shuffle} with the other two shuffling schemes outperforms SGD on training data.

% \section{Conclusion}

% In this paper, we have advanced the understanding of shuffling-type gradient algorithms under non-uniform smoothness assumptions and improved these algorithms with specific learning rates. We provided counterexamples to illustrate the existence and divergence of non-Lipschitz functions for (P). Our results show that these algorithms converge efficiently under a general bounded variance assumption. Additionally, we established robust convergence rates for both strongly convex and non-strongly convex cases. These convergence results outperform SGD also when Lipschitz smoothness is violated, which is demonstrated by numerical experiments. Future research can build on these findings to explore further generalizations and applications in various optimization contexts.

\section{Conclusion}
We revisited the convergence of shuffling-type gradient algorithms under relaxed smoothness assumptions, 
establishing their convergence for nonconvex, strongly convex, and non-strongly convex settings. 
By introducing a more general smoothness condition, we demonstrated that these methods achieve 
competitive convergence rates without requiring Lipschitz smoothness, extending their applicability 
to a broader range of optimization problems. Our analysis covers both random reshuffling and arbitrary 
shuffling schemes, showing that properly chosen step sizes can ensure efficient convergence in both cases. 
Numerical experiments further validate our theoretical findings, demonstrating that shuffling-type methods 
outperform SGD in non-Lipschitz scenarios. These results provide a foundation for future work on 
shuffling-based optimization, including its integration with variance reduction techniques, possible tighter bounds in certain situations and its application 
to large-scale machine learning problems.

%\newpage

\section*{Impact Statement}

This paper presents work whose goal is to advance the field
of Machine Learning. There are many potential societal
consequences of our work, none which we feel must be
specifically highlighted here.

\section*{Acknowledgments}
This work was partially supported by NSF IIS 2347592, 2348169, DBI 2405416, CCF 2348306, CNS 2347617.

\bibliography{icml_2025}


%%%%%%%%%%%%%%%%%%%%%%%%%%%%%%%%%%%%%%%%%%%%%%%%%%%%%%%%%%%%

\newpage
\onecolumn
\appendix

\section{Appendix / supplemental material}

\subsection{Nonconvex Case Analysis}

In this section we prove the theorems in section \ref{section:nonconvex}.

\subsubsection{Lemmas}
\label{appendix:lemmas}

In this part we use notations as defined in Theorem \ref{theorem:originalassumption}, for completeness we repeat them here:

$$H:=\frac{4\Delta_1}{\delta}, G:=\sup\{u\geq 0|u^2\leq 2\ell(2u)\cdot H\}<\infty,$$ $$G':=\sqrt{2(1+nA)}G+\sqrt{2n}\sigma, L:=\ell(2G').$$

We first state some lemmas that are useful in our proof. The following lemma is a natural corollary of Definition \ref{def:gs}, by the fact that $\ell$ is non-decreasing.

% We try to prove the following Lemma which is essential for our proof of convergence in reshuffling algorithms with $l$-smooth assumption. Basically it states that, with parameters chosen in theorems \ref{theorem:nonconvexconstant}, with probability at least $1-\delta/2$, we have Lipschitz smoothness along the training trajectory. Define $\Delta_1:=F(w^{(1)}_0)-F^*\geq 0.$

% \begin{Lemma}
% \label{Lemma:lsmooth}
%     For any $0<\delta<1$, we denote $H:=\frac{4\Delta_1}{\delta}$, $G:=\sup\{u\geq 0|u^2\leq 2\ell(2u)\cdot H\}<\infty$, $G':=G+\sigma$, $L:=\ell(2G')$, $r:=G/L$. Implement Algorithm \ref{alg:shuffle} with stepsize $$\eta_t\leq \frac{G}{2G'L}, \; \sum_{t=1}^T \eta_t^3 \leq \frac{n\Delta_1}{L^2\sigma^2},$$ and initialization $w_0^{(1)}$ such that $F(w_0^{(1)})-F^*\le H$. Then with probability at least $1-\delta/2$, the variables in Algorithm \ref{alg:shuffle} satisfy the following Lipschitz properties for any $t\in[T]$, $i,j\in[n]$,
%     \begin{align}
%         \|\nabla F(w^{(t)}_0)-\nabla F(w^{(t+1)}_0)\|&\leq L\|w^{(t)}_0-w^{(t+1)}_0\|; \label{eq:FLip}\\
%         \|\nabla f(w^{(t)}_0;i)-\nabla f(w^{(t)}_j;i)\|&\leq L\|w^{(t)}_0-w^{(t)}_j\|.\label{eq:fi_Lip}
%     \end{align}
% \end{Lemma}

% In order to prove this Lemma, we first restate the definition of $\ell$-smoothness with our specific function and parameter choices.

\begin{Lemma}
    \label{Lemma:rlsmooth}
    If $F$ is $\ell$-smooth, for any $w\in$ $\mathrm{dom}(F)$ satisfying $\|\nabla F(w)\|\leq G$, we have $\mathcal{B}(w,G/\ell(2G))\subseteq$$\mathrm{dom}(F)$. For any $w_1, w_2\in \mathcal{B}(w,G/\ell(2G)))$, 
    $$\|\nabla F(w_1)-\nabla F(w_2)\|\leq \ell(2G)\|w_1-w_2\|,$$ 
    $$ F(w_1)\leq F(w_2)+\langle \nabla F(w_2), w_1-w_2\rangle +\frac{\ell(2G)}{2}\|w_1-w_2\|^2.$$
\end{Lemma}

The following lemma gives relationship between $\|\nabla f(w;i)\|$ and $\|\nabla F(w)\|$.

\begin{Lemma}
   \label{Lemma:G'}
   If Assumption \ref{assumption:varianceexpectation} is true, we have 
   $$\|\nabla f(w;i)\|\leq \sqrt{2(1+nA)}\|\nabla F(w)\|+\sqrt{2n}\sigma.$$
\end{Lemma}

\begin{proof}
    From Assumption \ref{assumption:varianceexpectation} we have that 
    \begin{align*}
    \|\nabla f(w;i)\|^2 & \leq 2\|\nabla f(w;i)-\nabla F(w)\|^2+2\|\nabla F(w)\|^2 \\
    & \leq 2\sum_{i=1}^n\|\nabla f(w;i)-\nabla F(w)\|^2+2\|\nabla F(w)\|^2 \\
    & \leq 2nA\|\nabla F(w)\|^2+2n\sigma^2+2\|\nabla F(w)\|^2 \\
    & = 2(1+nA)\|\nabla F(w)\|^2+2n\sigma^2.
    \end{align*}
    Taking square root on both sides and notice that $\|\nabla F(w)\|\geq 0$, $\sigma>0$ we have the conclusion.
\end{proof}

According to Lemma \ref{Lemma:G'}, for $w$ such that $\|\nabla F(w)\|\leq G$ is true, we have $$\|\nabla f(w;i)\|\leq \sqrt{2(1+nA)}G+\sqrt{2n}\sigma=G'$$ holds for all $i\in[n]$.

In our proof, we want that with high probability, $L$-Lipschitz smoothness in Lemma \ref{Lemma:rlsmooth}, for both $F(w)$ and $f(w;i)$, between $w^{(t)}_0$ and $w^{(t)}_j$ is true, for $t\in[T]$, $i, j\in[n]$. For that purpose, we can prove the following inequalities with high probability, for $t\in[T]$:
\begin{align}
    \|\nabla F(w^{(t)}_0)\|& \leq G; \notag \\
    \|w^{(t)}_0-w^{(t)}_n\|& \leq G'/\ell(G+G');\notag \\
    \|\nabla f(w^{(t)}_0;i)\| &\leq G', \forall i\in[n]; \notag \\
    \|w^{(t)}_0-w^{(t)}_j\|&\leq G'/\ell(2G'),\forall j\in[n].\label{eq:tobeproven}
\end{align}

By Lemma \ref{Lemma:G'} we know that the third inequality in (\ref{eq:tobeproven}) holds if the first inequality is true. Noticing that $\ell(G+G')\leq \ell(2G')$, it suffices to prove that, for $t\in[T]$,
\begin{equation} 
\label{eq:easiertoprove}
\|\nabla F(w^{(t)}_0)\| \leq G, \|w^{(t)}_0-w^{(t)}_j\|\leq G'/\ell(2G'), \forall j\in[n].
\end{equation}

For the first inequality, it can be hard to bound the gradient norm directly. The following lemma states the connection between gradient norm and function value of an $\ell$-smooth function. 

\begin{Lemma} (Lemma 3.5 in \citet{li2023convex})
\label{Lemma:gradienttodiff}
    If $F$ is $\ell$-smooth, then $$\|\nabla F(w)\|^2\leq 2\ell(2\|\nabla F(w)\|)\cdot (F(w)-F^*)$$ for any $w\in$$\mathrm{dom}(F)$.
\end{Lemma}

Since $\ell$ is sub-quadratic, with Lemma \ref{Lemma:gradienttodiff} we can bound the gradient norm by bounding the difference between the function value and the optimal value. To ease the proof, let us define the following stopping time:
\begin{align*}
    \tau := \min\{t|F(w^{(t)}_0)-F^*>H\}\wedge (T+1).
\end{align*}

For $t<\tau$, we have $\|\nabla F(w^{(t)}_0)\|\leq G$ based on the definition of $\tau$ and Lemma \ref{Lemma:gradienttodiff}, so the first inequality in (\ref{eq:easiertoprove}) is satisfied. The following lemma proves that the other inequality in (\ref{eq:easiertoprove}) is true for $t<\tau$ as well, therefore guarantees the Lipschitz smoothness before $\tau$.

\begin{Lemma}
\label{Lemma:epsilon}
For $t<\tau$, $\eta_t \leq \frac{1}{2L}$, we have for all $k\in[n]$ and $t\in[T]$, $\|w^{(t)}_0-w^{(t)}_k\|^2\leq G'/\ell(2G').$
\end{Lemma}

\begin{proof}
    We use induction to prove that 
\[
w_j^{(t)} \in \mathcal{B}(w_0^{(t)}, \frac{G'}{\ell(2G')}), j = 0, 1, \ldots, n.
\]
First of all, this claim is true for \( j = 0 \). Now suppose the claim is true for $j\leq k-1$, i.e., 
\[
\|w_0^{(t)} - w_j^{(t)}\| \leq \frac{G'}{\ell(2G')}, j = 0, 1, \ldots, k-1,
\]
we try to prove it for \( w_k^{(t)} \). From Lemma \ref{Lemma:rlsmooth}, we have Lipschitz smoothness, for all $f(w;i)$, between $w^{(t)}_0$ and $w^{(t)}_j$, if $j\leq k-1$.

Since we have  
\[
\|\nabla f(w_0^{(t)};i)\| \leq G', \, \forall i \in [n],
\]
for any \( i \in [n] \) and \( j \in [k-1] \) we have
\[
\|\nabla f(w_j^{(t)};i)\| \leq \|\nabla f(w_0^{(t)};i)\| + \|\nabla f(w_j^{(t)};i) - \nabla f(w_0^{(t)};i)\| \leq G' + L \|w_j^{(t)} - w_0^{(t)}\| \leq 2G'.
\]
Hence, by the algorithm design we have
\[
\|w_k^{(t)} - w_0^{(t)}\| = \left\|\sum_{j=0}^{k-1} \frac{\eta_t}{n} \nabla f(w_j^{(t)}; \pi_j^{(t)})\right\| \leq \sum_{j=0}^{k-1} \frac{\eta_t}{n} \|\nabla f(w_j^{(t)}; \pi_j^{(t)})\| \leq \sum_{j=0}^{k-1} \frac{2G' \eta_t}{n} \leq \frac{G'}{L} = \frac{G'}{\ell(2G')},
\]
where the third inequality uses \( k \leq n \) and \( \eta_t \leq \frac{1}{2L}.
\) By induction, the claim is true.
\end{proof}

% \begin{Lemma}
% \label{Lemma:epsilon}
%     For $t<\tau$, $\eta_t\leq \frac{r}{G'+B}$, we have for all $k\in[n]$ and $t\in[T]$, $\|\epsilon^{(t)}_k\|^2\leq L^2\eta_t^2\cdot (2G^2+3\sigma^2).$
% \end{Lemma}

% \begin{proof}
%     \begin{align*}
%         & \|\epsilon^{(t)}_k\|^2 \\
%         = \; & \Big\|\frac{1}{n}\sum_{j=0}^{k-1} (\nabla f(w^{(t)}_j;\pi^{(t)}_j)-\nabla f(w^{(t)}_0;\pi^{(t)}_j))\Big\|^2\\
%         \leq \; & \frac{1}{n}\sum_{j=0}^{k-1}\|\nabla f(w^{(t)}_j;\pi^{(t)}_j)-\nabla f(w^{(t)}_0;\pi^{(t)}_j)\|^2 \\
%         \leq \; & \frac{L^2}{n}\sum_{j=0}^{n-1}\|w^{(t)}_j-w^{(t)}_0\|^2\\
%         \leq \; & L^2\eta_t^2\cdot (2G^2+3\sigma^2).
%     \end{align*}
%     Here the first inequality uses Cauchy-Schwarz inequality, the second uses Lemma \ref{Lemma:rlsmooth} and the last uses Lemma \ref{Lemma:sumofwdiff}.
% \end{proof}

% Therefore, we have for any $t\in[T]$ and $k\in[n]$,
% \begin{align}
%     \|w^{(t)}_0-w^{(t)}_k\|=&\frac{\eta_t}{n}\Big\|\sum_{j=0}^{k-1}\nabla f(w^{(t)}_{j};\pi^{(t)}_{j+1}
% )\Big\|\nonumber\\
%     =&\eta_t \Big\|\frac{1}{n}\sum_{j=0}^{k-1}\nabla f(w^{(t)}_{0};\pi^{(t)}_{j+1})+\epsilon_k^{(t)}\Big\|\nonumber\\
%     \le&\eta_t(G'+L\eta_t\sqrt{2G^2+3\sigma^2})\\
%     \le & \eta_t(G'+\frac{\sqrt{2G^2+3\sigma^2}}{2}).
% \end{align}

% The last inequality is from $\eta_t\leq \frac{1}{2L}$. By $\eta_t \leq \frac{G}{L\sqrt{2G^2+3\sigma^2}}$ and $\eta_t\leq \frac{G}{2G'L}$ we have $$\|w^{(t)}_0-w^{(t)}_k\|\leq \frac{G}{2L}+\frac{G}{2L}=\frac{G}{L}=r.$$ 
Therefore, we have the desired Lipschitz smoothness property in Lemma \ref{Lemma:rlsmooth} for $t<\tau$. Our next target is to bound $\mathbb{P}(\tau\leq T).$

To simplify the notations, let us define
$$\epsilon^{(t)}_k:=\frac{1}{n}\sum_{j=0}^{k-1}(\nabla f(w^{(t)}_j;\pi^{(t)}_{j+1})-\nabla f(w^{(t)}_0;\pi^{(t)}_{j+1}))$$ as the average of differences between the gradients at the start of iteration $t$ and the actual gradients we used until step $j$ in the $t$-th outer iteration. It is worth mentioning that the actual step in $t$-th outer iteration is $-\eta_t[\nabla F(w^{(t)}_0)+\epsilon^{(t)}_n]$.

Now we bound the probability of event $\{\tau\leq T\}$ by bounding the expectation of function value at the stopping time.

\begin{Lemma}
\label{Lemma:Ftau}
    With parameters chosen in Theorem \ref{theorem:originalassumption}, we have $$\mathbb{E}[F(w^{(\tau)}_0)-F^*]\leq 2\Delta_1.$$
    % where $j=\min\{j|F(w^{(\tau)}_j)-F^*>H\}$ if $\tau<T$ and $j=0$ otherwise.
\end{Lemma}

\begin{proof}
    For any $t<\tau$,
    \begin{align}
        & F(w^{(t+1)}_0)-F(w^{(t)}_0) \notag \\
        \leq & \langle \nabla F(w^{(t)}_0), w^{(t+1)}_0-w^{(t)}_0\rangle + \frac{L}{2}\|w^{(t+1)}_0-w^{(t)}_0\|^2 \notag \\
        = & -\eta_t\langle \nabla F(w^{(t)}_0), \nabla F(w^{(t)}_0)+\epsilon^{(t)}_n\rangle + \frac{L\eta_t^2}{2}\|\nabla F(w^{(t)}_0)+\epsilon^{(t)}_n\|^2 \notag \\
        = & -\frac{\eta_t}{2}(\|\nabla F(w^{(t)}_0)\|^2+\|\nabla F(w^{(t)}_0)+\epsilon^{(t)}_n\|^2-\|\epsilon^{(t)}_n\|^2)+\frac{L\eta_t^2}{2}\|\nabla F(w^{(t)}_0)+\epsilon^{(t)}_n\|^2 \notag \\
        \leq & -\frac{\eta_t}{2}\|\nabla F(w^{(t)}_0)\|^2+\frac{\eta_t}{2}\|\epsilon^{(t)}_n\|^2 \notag \\
        \leq & -\frac{\eta_t}{2}\|\nabla F(w^{(t)}_0)\|^2+\frac{\eta_t L^2}{2n}\sum_{k=0}^{n-1} \|w^{(t)}_k-w^{(t)}_0\|^2. \label{ineq:deltaF}
    \end{align}
    Here the first and last inequalities are from Lemma \ref{Lemma:rlsmooth} and the second is because $\eta_t\leq\frac{1}{2L}$. Taking summation from $t=1$ to $t=\tau-1$ and taking expectation we have
    \begin{equation}
    \mathbb{E}[F(w^{(\tau)}_0)-F^*]\leq \Delta_1-\mathbb{E}[\sum_{t=1}^{\tau-1}\frac{\eta_t}{2}\|\nabla F(w^{(t)}_0)\|^2]+\mathbb{E}[\sum_{t=1}^{\tau-1}\frac{\eta_t L^2}{2n}\sum_{k=0}^{n-1}\|w^{(t)}_k-w^{(t)}_0\|^2]. \label{ineq:F}
    \end{equation}

Now let us get a bound for the last term on the right hand side. For any $t\in[T]$, $k\in[n]$, from Algorithm \ref{alg:shuffle} and Cauchy-Schwarz inequality we have 
\begin{align}
    \|w^{(t)}_k-w^{(t)}_0\|^2 = &\frac{k^2\eta_t^2}{n^2}\Big\|\frac{1}{k}\sum_{j=0}^{k-1} \nabla f(w^{(t)}_j;\pi^{(t)}_{j+1})\Big\|^2 \notag \\
    \leq & \frac{3k^2\eta_t^2}{n^2}\|\frac{1}{k}\sum_{j=0}^{k-1}(\nabla f(w^{(t)}_0;\pi^{(t)}_{j+1})-\nabla F(w^{(t)}_0))\|^2+\frac{3k^2\eta_t^2}{n^2}\|\nabla F(w^{(t)}_0)\|^2 \notag \\
    & +\frac{3k\eta_t^2}{n^2}\sum_{j=0}^{k-1}\|\nabla f(w^{(t)}_j;\pi^{(t)}_{j+1})-\nabla f(w^{(t)}_0;\pi^{(t)}_{j+1})\|^2. \notag
\end{align}
    Let us denote the 3 terms on the RHS as $A_1(t, k), A_2(t, k)$ and $A_3(t, k)$, i.e. $\|w^{(t)}_k-w^{(t)}_0\|^2\leq A_1(t,k)+A_2(t,k)+A_3(t,k)$. Since we are interested in $\mathbb{E}[\sum_{t=1}^{\tau-1}\frac{\eta_t L^2}{2n}\sum_{k=0}^{n-1}\|w^{(t)}_k-w^{(t)}_0\|^2]$, we need to bound $\mathbb{E}[\sum_{t=0}^{\tau-1}\frac{\eta_t L^2}{2n}\sum_{k=1}^{n-1}A_i(t, k)]$ for $i=1,2,3.$

    For $A_1(t, k)$, since $\pi^{(t)}$ is randomly chosen, let $\mathcal{F}_t:=\sigma(\pi^{(1)},\cdots,\pi^{(t)})$ be the $\sigma$-algebra generated in Algorithm \ref{alg:shuffle}, for $t\in[T]$ we have 
    \begin{align*}
        \mathbb{E}[\frac{\eta_t L^2}{2n}\sum_{k=0}^{n-1}A_1(t, k)|\mathcal{F}_{t-1}] &= \frac{\eta_t L^2}{2n}\sum_{k=0}^{n-1}\frac{3k^2\eta_t^2}{n^2} \mathbb{E}[\|\frac{1}{k}\sum_{j=0}^{k-1}\nabla f(w^{(t)}_0;\pi^{(t)}_{j+1})-\nabla F(w^{(t)}_0)\|^2|\mathcal{F}_{t-1}] \\
        &= \frac{\eta_t L^2}{2n}\sum_{k=0}^{n-1}\frac{3k^2\eta_t^2}{n^2}\frac{n-k}{k(n-1)}\frac{1}{n}\sum_{i=0}^{n-1}\|\nabla f(w^{(t)}_0;i+1)-\nabla F(w^{(t)}_0)\|^2 \\
        &\leq \frac{\eta_t L^2}{2n}\sum_{k=0}^{n-1} \frac{3\eta_t^2 k(n-k)}{n^2(n-1)}(A\|\nabla F(w^{(t)}_0)\|^2+\sigma^2)\\
        &\leq \frac{\eta_t^3 L^2}{2n}(A\|\nabla F(w^{(t)}_0)\|^2+\sigma^2).
    \end{align*} 
    Here the second equation comes from variance of randomized reshuffling variables, (Lemma 1 in \citet{mishchenko2020random}); the first inequality is from assumption \ref{assumption:varianceexpectation}; the last inequality is because $\sum_{k=0}^{n-1} k(n-k)=\frac{(n-1)n(n+1)}{6}\leq \frac{n^2(n-1)}{3}$.
     
    Let $\{Z_t\}_{t\leq T}$ be a sequence such that $Z_1=0$ and for any $t\in[2,T]$, $$Z_{t}-Z_{t-1} = -\frac{\eta_t^3 L^2}{2n}(A\|\nabla F(w^{(t)}_0)\|^2+\sigma^2)+\frac{\eta_t L^2}{2n}\sum_{k=0}^{n-1}A_1(t-1,k).$$ We know $\{Z_t\}$ is a supermartingale. Since $\tau$ is a bounded stopping time, by optional stopping theorem, we have $\mathbb{E}[Z_{\tau}]\leq \mathbb{E}[Z_1]$, which leads to 
    $$\mathbb{E}[\sum_{t=1}^{\tau-1}\frac{\eta_t L^2}{2n}\sum_{k=0}^{n-1}A_1(t,k)]\leq \mathbb{E}[\sum_{t=1}^{\tau-1} \frac{\eta_t^3 L^2}{2n}(A\|\nabla F(w^{(t)}_0)\|^2+\sigma^2)].$$

    For $A_2(t,k)$, for any $t\in[T]$, taking summation over $k$ we have $\sum_{k=0}^{n-1}A_2(t,k)\leq n\eta_t^2\|\nabla F(w^{(t)}_0)\|^2$, therefore 
    $$\mathbb{E}[\sum_{t=1}^{\tau-1}\frac{\eta_t L^2}{2n}\sum_{k=0}^{n-1}A_2(t,k)]\leq \mathbb{E}[\sum_{t=1}^{\tau-1}\frac{\eta_t^3 L^2}{2}\|\nabla F(w^{(t)}_0)\|^2].$$

    For $A_3(t,k)$, for any $t<\tau$, by Lemma \ref{Lemma:rlsmooth} we have $$A_3(t,k)\leq \frac{3kL^2\eta_t^2}{n^2}\sum_{j=0}^{n-1}\|w^{(t)}_j-w^{(t)}_0\|^2.$$ Taking summation over $k$, taking expectation we have
    $$\mathbb{E}[\sum_{t=1}^{\tau-1}\frac{\eta_t L^2}{2n}\sum_{k=0}^{n-1}A_3(t,k)]\leq \mathbb{E}[\sum_{t=1}^{\tau-1}\frac{3\eta_t^3 L^4}{4n}\sum_{j=0}^{n-1}\|w^{(t)}_j-w^{(t)}_0\|^2].$$

    Now putting these together, we have 
    \begin{align*}
    \mathbb{E}[\sum_{t=1}^{\tau-1}\frac{\eta_t L^2}{2n}\sum_{j=0}^{n-1} \|w^{(t)}_j-w^{(t)}_0\|^2]\leq & \mathbb{E}[\sum_{t=1}^{\tau-1}\frac{\eta_t^3 L^2}{2n}(A\|\nabla F(w^{(t)}_0)\|^2+\sigma^2)]+\mathbb{E}[\sum_{t=1}^{\tau-1}\frac{\eta_t^3 L^2}{2}\|\nabla F(w^{(t)}_0)\|^2]\\
    & +[\sum_{t=1}^{\tau-1}\frac{3\eta_t^3 L^4}{4n}\sum_{j=0}^{n-1}\|w^{(t)}_j-w^{(t)}_0\|^2].
    \end{align*}

    Since $\eta_t\leq\frac{1}{2L}<\frac{1}{\sqrt{3}L}$ we have $\frac{3\eta_t^3 L^4}{4n}\leq \frac{\eta_t L^2}{4n}$, rearranging the terms we have 
    \begin{equation}
    \mathbb{E}[\sum_{t=1}^{\tau-1}\sum_{j=0}^{n-1} \|w^{(t)}_j-w^{(t)}_0\|^2]\leq \mathbb{E}[\sum_{t=1}^{\tau-1}2\eta_t^2\sigma^2]+2n\mathbb{E}[\sum_{t=1}^{\tau-1}\eta_t^2(\frac{A}{n}+1)\|\nabla F(w^{(t)}_0)\|^2]. \label{ineq:wdiff}
    \end{equation}

    Put this into (\ref{ineq:F}) we have, 
    \begin{align*}
        & \mathbb{E}[F(w^{(\tau)}_0)-F^*] \\
        \leq & \Delta_1+\mathbb{E}\Big[\sum_{t=1}^{\tau-1} \Big(-\frac{\eta_t}{2}\|\nabla F(w^{(t)}_0)\|^2+\frac{\eta_t L^2}{2n}\sum_{j=0}^{n-1} \|w^{(t)}_j-w^{(t)}_0\|^2\Big)\Big] \\
        \overset{(\ref{ineq:wdiff})}{\leq} & \Delta_1+\mathbb{E}\Big[\sum_{t=1}^{\tau-1}\frac{L^2\sigma^2\eta_t^3}{n}-\sum_{t=1}^{\tau-1}\Big((\frac{\eta_t}{2}-(\frac{A}{n}+1)\eta_t^3 L^2)\|\nabla F(w^{(t)}_0)\|^2\Big)\Big] \\
        \leq & \Delta_1+\mathbb{E}\Big[\sum_{t=1}^{\tau-1} \frac{L^2\sigma^2\eta_t^3}{n}-\sum_{t=1}^{\tau-1}\Big(\frac{\eta_t}{4}\|\nabla F(w^{(t)}_0)\|^2\Big)\Big] \tag{\theequation}\stepcounter{equation}\label{eq:expectation} \\
        \leq & \Delta_1+\frac{L^2\sigma^2}{n}\sum_{t=1}^T \eta_t^3.
    \end{align*}
    Here the third inequality is from $\eta_t\leq \frac{1}{2L\sqrt{\frac{A}{n}+1}}$ and the last inequality is because $\eta_t>0$ and $\tau\leq T+1$.
    
    Since $\sum_{t=1}^T \eta_t^3 \leq\frac{n\Delta_1}{\sigma^2 L^2}$, we have 
    $\mathbb{E}[F(w^{(\tau)}_0)-F^*]\leq 2\Delta_1.$
\end{proof}

Now we can bound the probability that $\tau=T+1$.

\begin{Lemma}
\label{Lemma:poftau}
    With the parameters in Theorem \ref{theorem:originalassumption}, we have
    $$\mathbb{P}(\tau\leq T)\leq \delta/2.$$
\end{Lemma}

\begin{proof}

    From Lemma \ref{Lemma:Ftau} and the value of $H$ we have

    $$\mathbb{P}(\tau\leq T)\leq \mathbb{P}(F(w^{(\tau)}_0)-F^*>H)\leq \frac{\mathbb{E}[F(w^{(\tau)}_0)-F^*]}{H}\leq\frac{2\Delta_1}{H}=\frac{\delta}{2}.$$
\end{proof}

\subsubsection{Proof for Theorems in Nonconvex cases}
\label{appendix:proofsfornonconvex}

\paragraph{Proof for Theorem \ref{theorem:originalassumption}}

\begin{proof}
    From (\ref{eq:expectation}) we have
    \begin{equation}
        \mathbb{E}[F(w^{\tau}_0)-F^*]+\mathbb{E}\Big[\sum_{t=1}^{\tau-1}\frac{\eta_t}{4}\|\nabla F(w^{(t)}_0)\|^2\Big]\leq \Delta_1+\frac{L^2\sigma^2}{n}\sum_{t=1}^T \eta_t^3\leq 2\Delta_1.
    \end{equation}

    Therefore, since $\delta\leq 1$ we have 
    \begin{align*}
        \frac{8\Delta_1}{\eta_T} & \geq \mathbb{E}\Big[\sum_{t=1}^{\tau-1}\|\nabla F(w^{(t)}_0)\|^2\Big] \\
        & \geq \mathbb{P}(\tau=T+1)\mathbb{E}\Big[\sum_{t=1}^T \|\nabla F(w^{(t)}_0)\|^2|\tau=T+1] \\
        & \geq \frac{1}{2}\mathbb{E}\Big[\sum_{t=1}^T\|\nabla F(w^{(t)}_0)\|^2|\tau=T+1\Big].
    \end{align*}

    By Markov's inequality and our choice of $T$, we have
    $$\mathbb{P}\Big(\frac{1}{T}\sum_{t=1}^T\|\nabla F(w^{(t)}_0)\|^2>\epsilon^2|\tau=T+1\Big)\leq \frac{16\Delta_1}{\eta_T T\epsilon^2}\leq \frac{\delta}{2}.$$
    
    From Lemma (\ref{Lemma:poftau}) we have $\mathbb{P}(\tau\leq T)\leq \frac{\delta}{2}$. Therefore, 
    \begin{align*}
        & \mathbb{P}\Big(\{\frac{1}{T}\sum_{t=1}^T\|\nabla F(w^{(t)}_0)\|^2>\epsilon^2\}\cup\{\tau\leq T\}\Big)\\
        \leq &\mathbb{P}(\tau\leq T)+\mathbb{P}\Big(\frac{1}{T}\sum_{t=1}^T\|\nabla F(w^{(t)}_0)\|^2>\epsilon^2|\tau=T+1)\\ 
        \leq &\frac{\delta}{2}+\frac{\delta}{2}=\delta.
    \end{align*}
    Since $L=\ell(2G')=\Omega(G'^p)=\Omega(n^{\frac{p}{2}})$, with $\eta=\mathcal{O}(\sqrt[3]{\frac{n^{1-p}}{T}})$ and $T=\mathcal{O}(\frac{n^{\frac{p-1}{2}}}{\epsilon^3})$ we have the complexity.
\end{proof}

% \paragraph{Proof of Corollary \ref{theorem:subgaussiannonconvexrr}}

% \begin{proof}

%     The only difference between this theorem and theorem \ref{theorem:nonconvexconstant} is that we do not always have $\|\nabla f(w;i)\|\leq G'$ even if $\|\nabla F(w)\|\leq G$. But since the event $\mathcal{F}:=\{\exists t\in[T], i\in[n], \text{ s.t. } \|\nabla f(w^{(t)}_0;i)-\nabla F(w^{(t)}_0)\|> \sigma$ for given $\sigma>0\}$ can be bounded from Assumption \ref{assumption:subgaussian}, we need to choose appropriate parameters so that this event happen with small probability and the total complexity does not increase much.

%     By setting $\sigma=R\sqrt{2\log(\frac{8nT}{\delta})}$, we have 
%     \begin{align*}
%         \mathbb{P}(\mathcal{F}) & = \mathbb{P}\Big(\cup_{t\in[T],i\in[n]} \|\nabla f(w^{(t)}_0;i)-\nabla F(w^{(t)}_0)\|>\sigma\Big)\\
%         & \leq \sum_{t\in[T],i\in[n]} \mathbb{P}\Big(\|\nabla f(w^{(t)}_0;i)-\nabla F(w^{(t)}_0)\|>\sigma\Big)\\
%         & \leq \sum_{t\in[T],i\in[n]} 2\exp(-\frac{\sigma^2}{2R^2}) \\
%         & \leq 2nT\exp(-\frac{\sigma^2}{2R^2})\leq \frac{\delta}{4}.
%     \end{align*}

%     Then if we change parameters so that $\mathbb{P}(\tau\leq T)\leq \frac{\delta}{4}$, we have the conclusion in the theorem.
% \end{proof}

% \paragraph{Proof for Corollary \ref{theorem:originalassumption}}

% \begin{proof}
%     From Assumption \ref{assumption:varianceexpectation} we have that 
%     \begin{align*}
%     \|\nabla f(w;i)\|^2 & \leq 2\|\nabla f(w;i)-\nabla F(w)\|^2+2\|\nabla F(w)\|^2 \\
%     & \leq 2\sum_{i=1}^n\|\nabla f(w;i)-\nabla F(w)\|^2+2\|\nabla F(w)\|^2 \\
%     & \leq 2nA\|\nabla F(w)\|^2+2n\sigma^2+2\|\nabla F(w)\|^2 \\
%     & = 2(1+nA)\|\nabla F(w)\|^2+2n\sigma^2.
%     \end{align*}
%     Therefore, for $w$ such that $\|\nabla F(w)\|\leq G$ is true, we have $$\|f(w;i)\|\leq \sqrt{2(1+nA)G^2+2n\sigma^2}\leq \sqrt{2(1+nA)}G+\sqrt{2n}\sigma$$ holds for all $i\in[n].$ That means, by the same proof we have for Theorem \ref{theorem:nonconvexconstant}, we have the same result except that we replace $G'$ with 
%     $$G'=\sqrt{2(1+nA)}G+\sqrt{2n}\sigma=\Theta(\sqrt{n}),$$
%     and replace $\sigma^2$ in the constraints with $AG^2+\sigma^2$. Since $L=\ell(2G')=\Theta(G'^p)=\Theta(n^{\frac{p}{2}})$, with $\eta=\mathcal{O}(\sqrt[3]{\frac{n^{1-p}}{T}})$ and $T=\mathcal{O}(\frac{n^{\frac{p-1}{2}}}{\epsilon^3})$ we have the complexity. 
% \end{proof}

The following lemma is useful in the proof of arbitrary scheme.

\begin{Lemma}(lemma 6 in \citep{nguyen2021unified})
    \label{Lemma:sumofwdiff}
    For $t<\tau$ and $0<\eta_t\leq \frac{1}{L\sqrt{3}}$, we have
    $$\sum_{j=0}^{n-1}\|w^{(t)}_j-w^{(t)}_0\|^2\leq n\eta_t^2[(3A+2)\|\nabla F(w^{(t)}_0)\|^2+3\sigma^2].$$
\end{Lemma}

\paragraph{Proof for Theorem \ref{theorem:nonconvexallscheme}}

\begin{proof}

    From inequality (\ref{ineq:deltaF}) we have for any $t<\tau$,
    \begin{align*}
        & F(w^{(t+1)}_0)-F(w^{(t)}_0) \\
        \leq & -\frac{\eta_t}{2}\|\nabla F(w^{(t)}_0)\|^2+\frac{\eta_t L^2}{2n}\sum_{k=0}^{n-1} \|w^{(t)}_k-w^{(t)}_0\|^2 \\
        \leq & -\frac{\eta_t}{2}\|\nabla F(w^{(t)}_0)\|^2+\frac{\eta_t^3 L^2[(3A+2)\|\nabla F(w^{(t)}_0)\|^2+3\sigma^2]}{2}\\
        \leq & -\frac{\eta_t}{4}\|\nabla F(w^{(t)}_0)\|^2+\frac{3\eta_t^3 L^2\sigma^2}{2},
    \end{align*}
    where the second inequality is from Lemma \ref{Lemma:sumofwdiff} and the last inequality is from $\eta_t\leq \frac{1}{L\sqrt{2(3A+2)}}$. Now taking summation of $t$ from $1$ to $\tau-1$ we have
    $$F(w^{(\tau)}_0)-F^*\leq F(w^{(\tau)}_0)-F^* + \sum_{t=1}^{\tau-1}\frac{\eta_t}{4}\|\nabla F(w^{(t)}_0)\|^2\leq \Delta_1+\frac{3L^2\sigma^2}{2}\sum_{t=1}^{\tau-1}\eta_t^3\leq 2\Delta_1,$$
    where the last inequality is because $\tau\leq T+1$ and the choice of $\eta_t$. Therefore we have $\tau=T+1$ since $H\geq 2\Delta_1$. On the other hand, we also have 
    \begin{align*}
        \frac{8\Delta_1}{\eta_T} & \geq \sum_{t=1}^{\tau-1}\|\nabla F(w^{(t)}_0)\|^2 \\
        & = \sum_{t=1}^T\|\nabla F(w^{(t)}_0)\|^2.
    \end{align*}
    Therefore, we have  
    $$\frac{1}{T}\sum_{t=1}^T\|\nabla F(w^{(t)}_0)\|^2\leq \frac{8\Delta_1}{T\eta_T}\leq \epsilon^2$$
    from our choice of $T$.
\end{proof}

\subsection{Strongly Convex Case Analysis}
\label{stronglyproof}

\begin{Lemma}
\label{Lemma:valueofH}
    If we let $H\geq\frac{3\sigma^2}{4\mu}\log(\frac{4}{\epsilon})+\Delta_1$ and $\eta_t=\eta$, we have $\tau\geq \frac{2}{\mu\eta}\log(\frac{4}{\epsilon})$.
\end{Lemma}

\begin{proof}
    From inequality (\ref{ineq:deltaF}) we have for $t<\tau$
    \begin{align*}
        F(w^{(t+1)}_0)-F(w^{(t)}_0) & \leq -\frac{\eta}{2}\|\nabla F(w^{(t)}_0)\|^2+\frac{\eta L^2}{2n}\sum_{k=0}^{n-1} \|w^{(t)}_k-w^{(t)}_0\|^2 \\
        & \leq -\frac{\eta}{2}\|\nabla F(w^{(t)}_0)\|^2+\frac{\eta^3 L^2[(3A+2)\|\nabla F(w^{(t)}_0)\|^2+3\sigma^2]}{2} \\
        & \leq \frac{3\eta\sigma^2}{8},
    \end{align*}
    where the last inequality is from $\eta\leq \frac{1}{L\sqrt{2(3A+2)}}\leq \frac{1}{2L}$. From the definition of $\tau$ we have 
    $$\tau\geq 1+\frac{8(H-\Delta_1)}{3\eta\sigma^2}\geq \frac{2}{\mu\eta}\log(\frac{4}{\epsilon}).$$
    
\end{proof}

\paragraph{Proof for Theorem \ref{theorem:stronglyconvexrr}}

\begin{proof} 

    From Lemma \ref{Lemma:poftau} and the parameter choices we have $\mathbb{P}(\tau\leq T)\leq \frac{\delta}{2}.$

    Now we try to bound $F(w^{(\tau)}_0)-F^*$. In the strongly convex case, for $t<\tau$ we have 
    \begin{align*}
        &F(w^{(t+1)}_0) \leq F(w^{(t)}_0)-\frac{\eta_t}{2}\|\nabla F(w^{(t)}_0)\|^2+\frac{L^2\eta_t}{2n}\sum_{j=0}^{n-1}\|w^{(t)}_j-w^{(t)}_0\|^2 \\
        & \leq F(w^{(t)}_0)-\frac{\mu\eta_t}{2} (F(w^{(t)}_0)-F^*)-\frac{\eta_t}{4}\|\nabla F(w^{(t)}_0)\|^2+\frac{L^2\eta_t}{2n}\sum_{j=0}^{n-1}\|w^{(t)}_j-w^{(t)}_0\|^2,
    \end{align*}
    here the first inequality is from (\ref{ineq:deltaF}) and the second one is from strongly convexity. We can rearrange the items and write the above inequality as 
    \begin{equation}
        F(w^{(t+1)}_0)-F^* \leq \Big(1-\frac{\mu\eta_t}{2}\Big)(F(w^{(t)}_0)-F^*)+\frac{L^2\sigma^2\eta_t^3}{n}+A(t),
    \end{equation}
    where $A(t)$ is defined as
    \begin{equation}
        A(t):= \frac{L^2\eta_t}{2n}\sum_{j=0}^{n-1}\|w^{(t)}_j-w^{(t)}_0\|^2-\frac{\eta_t}{4}\|\nabla F(w^{(t)}_0)\|^2-\frac{L^2\sigma^2\eta_t^3}{n}.
    \end{equation}
   Let $\eta_t=\eta:=\frac{4\log(\sqrt{n}T)}{\mu T}$, we want $1-\frac{\mu\eta}{2}>0$, therefore we need $\frac{T}{\log(\sqrt{n}T)}\geq 2$. Taking expectation and summation we have
    \begin{align}
        \mathbb{E}[F(w^{(\tau)}_0)-F^*] & \leq \mathbb{E}[(1-\frac{\mu\eta}{2})^{\tau-1}\Delta_1]+\frac{2L^2\sigma^2\eta^2}{n\mu}[1-(1-\frac{\mu\eta}{2})^{\tau-1}] \notag\\
        &+\mathbb{E}[\sum_{t=1}^{\tau-1}(1-\frac{\mu\eta}{2})^{\tau-1-t}A(t)]\notag \\
        &\leq \Delta_1\mathbb{E}[\exp(-\mu\eta \tau/2)] +\frac{2L^2\sigma^2\eta^2}{n\mu}+\mathbb{E}[\sum_{t=1}^{\tau-1}A(t)]\notag\\
        &\leq \frac{\delta\epsilon}{8}+\frac{1}{nT^2}\Big(\Delta_1+\frac{L^2\sigma^2\log^2(\sqrt{n}T)}{\mu^3}\Big)+\mathbb{E}[\sum_{t=1}^{\tau-1}A(t)],\notag
    \end{align}
    where the second inequality is from $1-x\leq \exp(-x)$ for $x\in(0,1)$ and the last inequality is from Lemma \ref{Lemma:valueofH}, $\mathbb{P}(\tau\leq T)\leq \delta/2$ and the value of $\eta$. Now if we look at the last item, we can notice from (\ref{ineq:wdiff}), by using $\eta\leq \frac{1}{L\sqrt{2(3A+2)}}\leq\frac{1}{2L\sqrt{\frac{A}{n}+1}}$, that we already have $$\mathbb{E}[\sum_{t=1}^{\tau-1}A(t)]\leq 0.$$

    Therefore, we have 
    \begin{align*}
        \frac{\delta\epsilon}{8}+\frac{1}{nT^2}\Big(\Delta_1+\frac{8L^2\sigma^2\log^2(\sqrt{n}T)}{\mu^3}\Big)&\geq \mathbb{E}[F(w^{(\tau)}_0)-F^*] \\
        & \geq \mathbb{P}(\tau=T+1)\mathbb{E}[F(w^{(T+1)}_0)-F^*|\tau=T+1] \\
        & \geq \frac{1}{2}\mathbb{E}[F(w^{(T+1)}_0)-F^*|\tau=T+1].
    \end{align*}

    \begin{align*}
    \mathbb{P}(F(w^{(T+1)}_0)-F^*>\epsilon|\tau=T+1)&\leq \frac{\mathbb{E}[F(w^{(T+1)}_0)-F^*|\tau=T+1]}{\epsilon} \\
    &\leq \frac{\delta}{4}+\frac{2}{\epsilon nT^2}\Big(\Delta_1+\frac{8L^2\sigma^2\log^2(\sqrt{n}T)}{\mu^3}\Big)\\
    & \leq \frac{\delta}{4}+\frac{\delta}{8}+\frac{\delta}{8}=\frac{\delta}{2},
    \end{align*}
    where the last line is from the constraint on T.

    % Now we need to confirm that the parameter choices do not conflict. If we have the $\ell$ function $\ell(x)=\Theta(x^p)$ for some constant $0\leq p<2$, since we choose $H=\Theta(\log(\frac{1}{\epsilon})$ in Lemma \ref{Lemma:valueofH}, we have $G=\Theta(\log^{\frac{1}{2-p}}(\frac{1}{\epsilon}))$ and $L=\Theta(\log^{\frac{p}{2-p}}(\frac{1}{\epsilon}))$. On the other hand, from the choice of $T$ and the relationship between $\eta$ and $T$ we have $\eta\approx \mathcal{O}(\frac{\delta\epsilon}{n}\log(\frac{1}{\epsilon}))$. Therefore, the requirements $\eta_t\leq \min \Big\{\frac{1}{2L}, \frac{G}{2G'L}, \frac{G}{L\sqrt{2G^2+3\sigma^2}}\Big\}$ are met for small enough $\epsilon$.
\end{proof}

\paragraph{Proof for Theorem \ref{theorem:stronglyconvexany}}

\begin{proof}

The algorithm starts from $w^{(1)}_0$ and we define $S=\{w|F(w)\leq F(w^{(1)}_0)\}.$ Since $F$ is strongly-convex, we have $S$ being compact. Therefore, we can define $G'=\max_w\{\|\nabla f(w;i)\||w\in S, i\in[n]\}<\infty.$ 

If we have $w^{(t)}_0\in S$ for all $t\in[T]$, we have $\|\nabla f(w^{(t)}_0;i)\|\leq G'$ for $t\in[T]$ and $i\in [n]$. On the other hand, by definition of $F$ we have $\|\nabla F(w^{(t)}_0)\|\leq G'$ for $t\in[T]$. Therefore, by Lemma \ref{Lemma:epsilon} we have Lipschitz smoothness between $w^{(t)}_0$ and $w^{(t)}_j$, for both $F(w)$ and $f(w;i)$, for $t\in[T]$, $i,j\in[n]$. The rest of the proof then follows the one in Lipschitz smoothness case (theorem 1 in \citet{nguyen2021unified}).

Now we prove that $w^{(t)}_0\in S$, for $t\in T.$ The statement is obviously true for $t=1$. Now for $t\in[2, T]$, assume that we already proved the conclusion for $1,\cdots, t-1$, we can use Lipschitz smoothness in the first $t-1$ iterations. Therefore, from theorem 1 in \citet{nguyen2021unified} we have that 
$$F(w^{(t)}_0)-F(w_*)\leq (1-\rho \eta)^{t-1}\Delta_1 +\frac{D\eta ^2}{\rho},$$
where $\rho=\frac{\mu}{3}$, $D= (\mu^2+L^2)\sigma^2_*.$ On the other hand, since $\eta \leq \frac{\Delta_1\rho^2}{D}$ we have 
$$(1-\rho \eta)^{t-1}\Delta_1 +\frac{D\eta ^2}{\rho}\leq (1-\rho \eta)\Delta_1 +\frac{D\eta^2}{\rho}\leq \Delta_1. $$
Therefore, we have $F(w^{(t)}_0)\leq F(w^{(1)}_0)$, which means $w^{(t)}_0\in S.$

\end{proof}

% \begin{proof}

% Similar to proof of Theorem \ref{theorem:nonconvexallscheme} we have $\tau=T+1$. 

% Similar to proof of Theorem \ref{theorem:stronglyconvexrr} we have
% \begin{equation}
%     F(w^{(t+1)}_0)-F^* \leq \Big(1-\frac{\mu\eta_t}{2}\Big)(F(w^{(t)}_0)-F^*)+\frac{3L^2\sigma^2\eta_t^3}{2}+A(t),
% \end{equation}

% and this time $A(t)$ is defined as
% \begin{equation}
%     A(t):= \frac{L^2\eta_t}{2n}\sum_{j=0}^{n-1}\|w^{(t)}_j-w^{(t)}_0\|^2-\frac{\eta_t}{4}\|\nabla F(w^{(t)}_0)\|^2-\frac{3L^2\sigma^2\eta_t^3}{2}.
% \end{equation}

% From Lemma \ref{Lemma:sumofwdiff} and $\eta\leq\frac{1}{\sqrt{2(3A+2)}}$ we have $A(t)\leq 0.$ 

% Let $\eta_t=\eta:=\frac{4\log(\sqrt{n}T)}{\mu T}$, similarly we have

% \begin{align*}
%     F(w^{(T+1)}_0)-F^* & \leq (1-\frac{\mu\eta}{2})^{T}\Delta_1+\frac{3L^2\sigma^2\eta^2}{2\mu}[1-(1-\frac{\mu\eta}{2})^{T}] \notag\\
%     &+\sum_{t=1}^{T}(1-\frac{\mu\eta}{2})^{T-t}A(t)\notag \\
%     &\leq \Delta_1\exp(-\mu\eta T/2) +\frac{3L^2\sigma^2\eta^2}{2\mu}+\sum_{t=1}^{T}A(t)\\
%     &\leq \frac{1}{nT^2}\Big(\Delta_1+\frac{3nL^2\sigma^2\log^2(\sqrt{n}T)}{2\mu^3}\Big).
% \end{align*}

% Therefore, we have $F(w^{(T+1)}_0)-F^*\leq \epsilon$ by the choice of $T$.

% end here

% \begin{equation}\label{Ls}
% \begin{split}
%   F(w_0^{(t+1)})&\le F(w_0^{(t)}) + \left<\nabla F(w_0^{(t)}),w_0^{(t+1)} - w_0^{(t)}\right> + \frac{L}{2}\|w_0^{(t+1)} - w_0^{(t)}\|^2\\
%   &=F(w_0^{(t)}) - \eta\left<\nabla F(w_0^{(t)}),\nabla F(w_0^{(t)}) + \epsilon_n^{(t)}\right> + \frac{L\eta^2}{2}\|\nabla F(w_0^{(t)}) + \epsilon_n^{(t)}\|^2\\
%   &=F(w_0^{(t)}) - \frac12\eta\left(\|\nabla F(w_0^{(t)})\|^2 +  \|\nabla F(w_0^{(t)}) + \epsilon_n^{(t)}\|^2 - \|\epsilon_n^{(t)}\|^2\right)\\
%   &+ \frac{L\eta^2}{2}\|\nabla F(w_0^{(t)}) + \epsilon_n^{(t)}\|^2\\
%   &\stackrel{\rm (a)}\le F(w_0^{(t)}) - \frac12\eta\left(\|\nabla F(w_0^{(t)})\|^2 +   2\|\nabla F(w_0^{(t)})\|^2 + 2\|\epsilon_n^{(t)}\|^2 + \|\epsilon_n^{(t)}\|^2\right)\\
%   & + L\eta^2\left(\|\nabla F(w_0^{(t)})\|^2 + \|\epsilon_n^{(t)}\|^2\right)\\
%   &=F(w_0^{(t)}) - \left(\frac{3\eta}{2} - L\eta^2\right)\|\nabla F(w_0^{(t)})\|^2 + \left(\frac{3\eta}2 +L\eta^2\right)\|\epsilon_n^{(t)}\|^2,
% \end{split}
% \end{equation}
% where the first equality uses step 6 in Algorithm \ref{alg:shuffle} and definition of $\epsilon_0^{(t)}$, the second equality uses the fact that $\left<a,b\right>=  \frac12\|a\|^2 + \frac12\|b\|^2 -\frac12\|a-b\|^2$ for any vectors $a$ and $b$, (a) uses $\|a+b\|^2\le 2a^2 + 2b^2$. Since $F$ is strongly convex, there exists unique optimizer $w_*$ and constant $\mu>0$ such that the following inequality holds:
% \begin{equation}\label{PL}
%   2\mu\left(F(w) - F(w_*)\right)\le \|\nabla F(w)\|^2.
% \end{equation}
% Since $\eta\leq\frac{1}{2L}$, $\frac{3\eta}{2} - L\eta^2>0$. Therefore, using \eqref{PL}, \eqref{Ls} can be further passed to
% \begin{equation}
%   \begin{split}
%     F(w_0^{(t+1)})&\le F(w_0^{(t)}) - \left(\frac{3\eta}{2} - L\eta^2\right)2\mu\left(F(w^{(t)}_0) - F(w_*)\right) + \left(\frac{3\eta}2 +L\eta^2\right)\|\epsilon_n^{(t)}\|^2\\
%     &\le  F(w_0^{(t)}) - \left(\frac{3\eta}{2} - L\eta^2\right)2\mu\left(F(w^{(t)}_0) - F(w_*)\right) + \left(\frac{3\eta}2 +L\eta^2\right)L^2 (2G^2+3\sigma^2)\eta^2.
%   \end{split}
% \end{equation}
% Rearranging the above inequality, we have
% \begin{equation*}
%   \begin{split}
% &F(w^{(t+1)}_0)-F(w_*)\\
% &\le (1- \underbrace{\left(\frac{3}{2} - L\eta\right)2\mu}_{=:\rho}\eta)\left(F(w^{(t)}_0) - F(w_*)\right) + \left(\frac{3\eta}2 +L\eta^2\right)L^2(2G^2+3\sigma^2)\eta^2.
%   \end{split}
% \end{equation*}

% Since $\rho\leq (\frac32-0)2\mu=3\mu$ and $\frac{\log T}{T}<\frac14$ for all $T>0$, by setting $\eta=\frac{\log T}{\mu T}$ and $\eta\leq \frac{1}{2L}$, we have $\eta\rho<1$ and $\rho>0$. Thus, with similar computation as in \eqref{eq:stronglyconvex}, it holds that
% \begin{equation}\label{l7}
%   F(w^{(T+1)}_0)-F(w_*)\le \left(F(w^{(1)}_0)-F(w_*)\right)e^{-\rho\eta T} + \left(\frac{3\eta}2 +L\eta^2\right)L^2(2G^2+3\sigma^2)\frac{\eta}{\rho}.
% \end{equation}
% Since $\rho\geq 2\mu$, \eqref{l7} becomes
% \begin{equation*}
% \begin{split}
%   &F(w^{(T+1)}_0)-F(w_*)\\
%   &\le \left(F(w^{(1)}_0)-F(w_*)\right)\frac{1}{T^2} + 2\eta L^2(2G^2+3\sigma^2)\frac{\eta}{2\mu}\\
%   &\le \left(F(w^{(1)}_0)-F(w_*)\right)\frac{1}{T^2} + \frac{L^2(2G^2+3\sigma^2)}{\mu^3}\left(\frac{\log T}{T}\right)^2.
%   \end{split}
% \end{equation*}

% \end{proof}

\subsection{Non-strongly Convex Case Analysis}
\label{convexproof}

\paragraph{Proof for theorem \ref{theorem:nonstrongly}}

\begin{proof}

From Lemma \ref{Lemma:poftau} we know $\mathbb{P}(\tau\leq T)<\frac{\delta}{2}.$

    For $t<\tau$, if $\eta_t=\eta$, from lemma 7 in \citep{nguyen2021unified} we have that
    \begin{equation}
    \label{ineq:convex}
        \|w^{(t+1)}_0-w_*\|^2\leq \|w^{(t)}_0-w_*\|^2-2\eta[F(w^{(t)}_0)-F^*]+\frac{2L\eta^3}{n^3}\sum_{i=1}^{n-1}\|\sum_{j=i}^{n-1} \nabla f(w_*;\pi^{(t)}_{j+1})\|^2,
    \end{equation}
    where $w_*$ is the optimal solution. If we denote $A(t):=\sum_{i=1}^{n-1}\|\sum_{j=i}^{n-1} \nabla f(w_*;\pi^{(t)}_{j+1})\|^2$ and let $\sigma_*:=\sqrt{\frac{1}{n}\sum_{i=1}^n\|\nabla f(w_*;i)\|^2}$, we have that for any $t\in[T]$
    \begin{align*}
    \mathbb{E}[A(t)]&=\sum_{i=0}^{n-1}(n-i)^2\mathbb{E}\left[\left\|\frac{1}{n-i}\sum_{j=i}^{n-1}\nabla f(w_*;\pi^{(t)}_{j+1}-\nabla F(w_*)\right\|^2\right] \\
    & = \sum_{i=0}^{n-1}\frac{(n-i)^2i}{n(n-i)(n-1)}\sum_{j=0}^{n-1}\|\nabla f(w_*;\pi^{(t)}_{j+1})\|^2\\
    & = \frac{n(n+1)\sigma_*^2}{6}.
    \end{align*}
    By optional stopping theorem we know that 
    $$\mathbb{E}\Big[\sum_{t=1}^{\tau-1}\Big(A(t)-\frac{n(n+1)\sigma_*^2}{6}\Big)\Big]=0.$$
    Taking summation from $t=0$ to $\tau-1$ for (\ref{ineq:convex}) and taking expectation we have
    \begin{align*}
        2\eta\mathbb{E}[\sum_{t=1}^{\tau-1}(F(w^{(t)}_0)-F^*)] &\leq \|w^{(1)}_0-w_*\|^2+\frac{2L\eta^3}{n^3}\mathbb{E}[\sum_{t=1}^{\tau-1}\frac{n(n+1)\sigma_*^2}{6}]\\
        &\leq \|w^{(1)}_0-w_*\|^2+\frac{2LT\eta^3\sigma_*^2}{3n},
    \end{align*}
    where the second inequality uses $\tau\leq T+1.$ Therefore, we have
    \begin{align*}
        \frac{1}{2\eta}\Big(\|w^{(1)}_0-w_*\|^2+\frac{2LT\eta^3\sigma_*^2}{3n}\Big)&\geq \mathbb{E}[\sum_{t=1}^{\tau-1}(F(w^{(t)}_0)-F^*)]\\
        &\geq \frac{1}{2}\mathbb{E}[\sum_{t=1}^{T}(F(w^{(t)}_0)-F^*)|\tau=T+1].
    \end{align*}
    
    If we define $\bar{w}_{T}=\frac{1}{T}\sum_{t=1}^{T}w^{(t)}_0$, from convexity we have 
    $$F(\bar{w}_T)-F^*\leq \frac{1}{T}\sum_{t=1}^T [F(w^{(t)}_0)-F^*].$$

    Consider the event $\mathcal{F}:=\{F(\bar{w}_T)-F^*>\epsilon\}$, we have 
    \begin{align*}
        \mathbb{P}(\mathcal{F}|\tau=T+1)&\leq \mathbb{P}(\frac{1}{T}\sum_{t=1}^T (F(w^{(t)}_0)-F^*)>\epsilon|\tau=T+1)\\
        &\leq \frac{\mathbb{E}\left[\sum_{t=1}^T (F(w^{(t)}_0)-F^*)|\tau=T+1\right]}{T\epsilon}\\
        &\leq \frac{1}{\eta T\epsilon} \Big(\|w^{(1)}_0-w_*\|^2+\frac{2LT\eta^3\sigma_*^2}{3n}\Big)\\
        & \leq \frac{2}{\eta T\epsilon}\|w^{(1)}_0-w_*\|^2\\
        & \leq \frac{\delta}{2},
    \end{align*}
    where the last two inequalities are from the choices of $\eta$ and $T$, separately. 
\end{proof}

\paragraph{Proof for Theorem \ref{theorem:nonstronglyany}}

\begin{proof}

Similar to Theorem \ref{theorem:stronglyconvexany}, if we have $w^{(t)}_0\in S$ for $t\in[T]$, we have the desired Lipschitz smoothness.

Now we prove the conclusion by trying to prove that $w^{(t)}_0\in S$ for $t\in[T]$. The statement is obviously true for $t=1$. Now for $t\in[2, T]$, assume that we already proved the conclusion for $1,\cdots, t-1$, we can use Lipschitz smoothness in the first $t-1$ iterations. Therefore, from \eqref{ineq:convex} we have
$$ \|w^{(t)}_0-w_*\|^2\leq \|w^{(t-1)}_0-w_*\|^2-2\eta[F(w^{(t-1)}_0)-F^*]+\frac{2L\eta^3}{n^3}\sum_{i=1}^{n-1}\|\sum_{j=i}^{n-1} \nabla f(w_*;\pi^{(t-1)}_{j+1})\|^2.$$

If $F(w^{(t-1)}_0)-F^*\leq \epsilon$, we have the desired conclusion.

If $F(w^{(t-1)}_0)-F^*> \epsilon$, since $\eta \leq \frac{1}{G'}\sqrt\frac{3\epsilon}{L}$, we have 
\begin{align*}
    \|w^{(t+1)}_0-w_*\|^2 & \leq \|w^{(t)}_0-w_*\|^2-2\eta \epsilon+\frac{2L\eta^3}{n^3}\sum_{i=1}^{n-1}(n-i)^2 G'^2\\
    &\leq \|w^{(t)}_0-w_*\|^2-2\eta \epsilon+\frac{2L\eta^3}{n^3}\frac{n^3}{3}G'^2 \\
    & \leq \|w^{(t)}_0-w_*\|^2.
\end{align*}
Therefore, if $F(w^{(t)}_0)-F^*\geq \epsilon $ for $t\in[T]$, we have $w^{(t)}_0\in S$ for $t\in[T].$ Taking summation we have that 
$$2\eta \sum_{t=1}^T[F(w^{(t)}_0)-F(w_*)]\leq \|w^{(1)}_0-w_*\|^2+\frac{2LG'^2\eta^3T}{3}.$$
Therefore we have 
$$\frac{1}{T} \sum_{t=1}^T[F(w^{(t)}_0)-F(w_*)]\leq \frac{1}{2\eta T}\left(\|w^{(1)}_0-w_*\|^2+\frac{2LG'^2\eta^3T}{3}\right)\leq \epsilon.$$
However, this contradict the assumption that $F(w^{(t)}_0)-F^*\geq \epsilon $ for $t\in[T]$. Therefore, there must be $t\in[T]$ such that $F(w^{(t)}_0)-F^*\leq \epsilon$.

\end{proof}

% \begin{proof}

%     Similar to proof of Theorem \ref{theorem:nonconvexallscheme} we have $\tau=T+1$. 

%     From \eqref{ineq:convex}, by Assumption \ref{assumption:varianceexpectation} we have
%     $$\|w^{(t+1)}_0-w_*\|^2\leq \|w^{(t)}_0-w_*\|^2-2\eta[F(w^{(t)}_0)-F^*]+\frac{2L\eta^3}{n^3}\sum_{i=1}^{n-1}(n-i)^2\sigma^2.$$

%     Taking summation from $t=0$ to $T-1$, notice that $\sum_{i=1}^{n-1}i^2\leq \frac{n^3}{2}$ we have 
%     $$2\eta\sum_{t=1}^{T}(F(w^{(t)}_0)-F^*)\leq \|w^{(1)}_0-w_*\|^2+LT\eta^3\sigma^2.$$
%     Therefore we have
%     $$F(\bar{w}_T)-F^*\leq \sum_{t=1}^{T}(F(w^{(t)}_0)-F^*)\leq \frac{1}{2\eta T}\Big(\|w^{(1)}_0-w_*\|^2+LT\eta^3\sigma^2\Big)\leq \epsilon.$$
% \end{proof}

% \subsection{Convex Case without variance assumptions}
% \label{appendix:final}

% \subsubsection{Strongly Convex Case}

% \paragraph{Proof for Theorem \ref{theorem:noassum1}}

% \begin{proof}

% The algorithm starts from $w^{(1)}_0$ and we define $S=\{w|F(w)\leq F(w^{(1)}_0)\}.$ Since $F$ is strongly-convex, we have $S$ being compact. Therefore, we can define $G'=\max_w\{\|\nabla f(w;i)\||w\in S, i\in[n]\}<\infty.$ 

% If we have $w^{(t)}_0\in S$ for all $t\in[T]$, we have $\|\nabla f(w^{(t)}_0;i)\|\leq G'$ for $t\in[T]$ and $i\in [n]$. On the other hand, by definition of $F$ we have $\|\nabla F(w^{(t)}_0)\|\leq G'$ for $t\in[T]$. Therefore, by Lemma \ref{Lemma:epsilon} we have Lipschitz smoothness between $w^{(t)}_0$ and $w^{(t)}_j$, for both $F(w)$ and $f(w;i)$, for $t\in[T]$, $i,j\in[n]$. The rest of the proof then follows the one in Lipschitz smoothness case (theorem 1 in \citet{nguyen2021unified}).

% Now we prove that $w^{(t)}_0\in S$, for $t\in T.$ The statement is obviously true for $t=1$. Now for $t\in[2, T]$, assume that we already proved the conclusion for $1,\cdots, t-1$, we can use Lipschitz smoothness in the first $t-1$ iterations. Therefore, from theorem 1 in \citet{nguyen2021unified} we have that 
% $$F(w^{(t)}_0)-F(w_*)\leq (1-\rho \eta)^{t-1}\Delta_1 +\frac{D\eta ^2}{\rho},$$
% where $\rho=\frac{\mu}{3}$, $D= (\mu^2+L^2)\sigma^2_*.$ On the other hand, since $\eta \leq \frac{\Delta_1\rho^2}{D}$ we have 
% $$(1-\rho \eta)^{t-1}\Delta_1 +\frac{D\eta ^2}{\rho}\leq (1-\rho \eta)\Delta_1 +\frac{D\eta^2}{\rho}\leq \Delta_1. $$
% Therefore, we have $F(w^{(t)}_0)\leq F(w^{(1)}_0)$, which means $w^{(t)}_0\in S.$

% \end{proof}

% \subsubsection{Non-Strongly Convex Case}

% \paragraph{Proof for Theorem \ref{theorem:noassum2}}

% \begin{proof}

% This time we define the compact set $S=\{w|\|w-w_*\|\leq \|w^{(1)}_0-w_*\|\}.$ Let us define $G'=\max_w\{\|\nabla f(w;i)\||w\in S, i\in[n]\}<\infty$ and $\Delta = \max_w\{F(w)-F(w_*)|w\in S, i\in[n]\}<\infty.$ If we have $w^{(t)}_0\in S$ for $t\in[T]$, once again we have the desired Lipschitz smoothness.

% Now we prove that $w^{(t)}_0\in S$ for $t\in[T].$ The statement is obviously true for $t=1$. Now for $t\in[2, T]$, assume that we already proved the conclusion for $1,\cdots, t-1$, we can use Lipschitz smoothness in the first $t-1$ iterations. Therefore, from \eqref{ineq:convex} we have
% $$ \|w^{(t)}_0-w_*\|^2\leq \|w^{(t-1)}_0-w_*\|^2-2\eta[F(w^{(t-1)}_0)-F^*]+\frac{2L\eta^3}{n^3}\sum_{i=1}^{n-1}\|\sum_{j=i}^{n-1} \nabla f(w_*;\pi^{(t-1)}_{j+1})\|^2.$$

% If $F(w^{(t-1)}_0)-F^*<\epsilon$, we have the desired conclusion.

% If $F(w^{(t-1)}_0)-F^*\geq \epsilon$, since $\eta \leq \frac{1}{G'}\sqrt\frac{3\epsilon}{L}$, we have 
% \begin{align*}
%     \|w^{(t+1)}_0-w_*\|^2 & \leq \|w^{(t)}_0-w_*\|^2-2\eta \epsilon+\frac{2L\eta^3}{n^3}\sum_{i=1}^{n-1}(n-i)^2 G'^2\\
%     &\leq \|w^{(t)}_0-w_*\|^2-2\eta \epsilon+\frac{2L\eta^3}{n^3}\frac{n^3}{3}G'^2 \\
%     & \leq \|w^{(t)}_0-w_*\|^2.
% \end{align*}
% Therefore, if $F(w^{(t)}_0)-F^*\geq \epsilon $ for $t\in[T]$, we have $w^{(t)}_0\in S$ for $t\in[T].$ Taking summation we have that 
% $$2\eta \sum_{t=1}^T[F(w^{(t)}_0)-F(w_*)]\leq \|w^{(1)}_0-w_*\|^2+\frac{2LG'^2\eta^3}{3}.$$
% Therefore we have 
% $$\frac{1}{T} \sum_{t=1}^T[F(w^{(t)}_0)-F(w_*)]\leq \frac{1}{2\eta T}\left(\|w^{(1)}_0-w_*\|^2+\frac{2LG'^2\eta^3}{3}\right)\leq \epsilon.$$
% Hence the conclusion holds.

% \end{proof}

% %%%%%%%%%%%%%%%%%%%%%%%%%%%%%%%%%%%%%%%%%%%%%%%%%%%%%%%%%%%%

% \newpage
% \section*{NeurIPS Paper Checklist}

% \begin{enumerate}

% \item {\bf Claims}
%     \item[] Question: Do the main claims made in the abstract and introduction accurately reflect the paper's contributions and scope?
%     \item[] Answer: \answerYes{} % Replace by \answerYes{}, \answerNo{}, or \answerNA{}.
%     \item[] Justification: The main claims in the abstract and introduction are all stated or theoretically proven in the paper.
%     \item[] Guidelines:
%     \begin{itemize}
%         \item The answer NA means that the abstract and introduction do not include the claims made in the paper.
%         \item The abstract and/or introduction should clearly state the claims made, including the contributions made in the paper and important assumptions and limitations. A No or NA answer to this question will not be perceived well by the reviewers. 
%         \item The claims made should match theoretical and experimental results, and reflect how much the results can be expected to generalize to other settings. 
%         \item It is fine to include aspirational goals as motivation as long as it is clear that these goals are not attained by the paper. 
%     \end{itemize}

% \item {\bf Limitations}
%     \item[] Question: Does the paper discuss the limitations of the work performed by the authors?
%     \item[] Answer: \answerYes{} % Replace by \answerYes{}, \answerNo{}, or \answerNA{}.
%     \item[] Justification: Assumptions \ref{assumption:variance} and \ref{assumption:subgaussian} are stronger than standard bounded variance assumption. We point that out in the paper, explain why we need the stronger assumptions and justify the usage by referring to other papers with the same assumptions.
%     \item[] Guidelines:
%     \begin{itemize}
%         \item The answer NA means that the paper has no limitation while the answer No means that the paper has limitations, but those are not discussed in the paper. 
%         \item The authors are encouraged to create a separate "Limitations" section in their paper.
%         \item The paper should point out any strong assumptions and how robust the results are to violations of these assumptions (e.g., independence assumptions, noiseless settings, model well-specification, asymptotic approximations only holding locally). The authors should reflect on how these assumptions might be violated in practice and what the implications would be.
%         \item The authors should reflect on the scope of the claims made, e.g., if the approach was only tested on a few datasets or with a few runs. In general, empirical results often depend on implicit assumptions, which should be articulated.
%         \item The authors should reflect on the factors that influence the performance of the approach. For example, a facial recognition algorithm may perform poorly when image resolution is low or images are taken in low lighting. Or a speech-to-text system might not be used reliably to provide closed captions for online lectures because it fails to handle technical jargon.
%         \item The authors should discuss the computational efficiency of the proposed algorithms and how they scale with dataset size.
%         \item If applicable, the authors should discuss possible limitations of their approach to address problems of privacy and fairness.
%         \item While the authors might fear that complete honesty about limitations might be used by reviewers as grounds for rejection, a worse outcome might be that reviewers discover limitations that aren't acknowledged in the paper. The authors should use their best judgment and recognize that individual actions in favor of transparency play an important role in developing norms that preserve the integrity of the community. Reviewers will be specifically instructed to not penalize honesty concerning limitations.
%     \end{itemize}

% \item {\bf Theory Assumptions and Proofs}
%     \item[] Question: For each theoretical result, does the paper provide the full set of assumptions and a complete (and correct) proof?
%     \item[] Answer: \answerYes{} % Replace by \answerYes{}, \answerNo{}, or \answerNA{}.
%     \item[] Justification: We give proof to all the theorems in the Appendix.
%     \item[] Guidelines:
%     \begin{itemize}
%         \item The answer NA means that the paper does not include theoretical results. 
%         \item All the theorems, formulas, and proofs in the paper should be numbered and cross-referenced.
%         \item All assumptions should be clearly stated or referenced in the statement of any theorems.
%         \item The proofs can either appear in the main paper or the supplemental material, but if they appear in the supplemental material, the authors are encouraged to provide a short proof sketch to provide intuition. 
%         \item Inversely, any informal proof provided in the core of the paper should be complemented by formal proofs provided in appendix or supplemental material.
%         \item Theorems and Lemmas that the proof relies upon should be properly referenced. 
%     \end{itemize}

%     \item {\bf Experimental Result Reproducibility}
%     \item[] Question: Does the paper fully disclose all the information needed to reproduce the main experimental results of the paper to the extent that it affects the main claims and/or conclusions of the paper (regardless of whether the code and data are provided or not)?
%     \item[] Answer: \answerYes{} % Replace by \answerYes{}, \answerNo{}, or \answerNA{}.
%     \item[] Justification: We provided all the details needed to reproduce the experimental results, including the objective functions, data generation process, hyperparameters, etc. 
%     \item[] Guidelines:
%     \begin{itemize}
%         \item The answer NA means that the paper does not include experiments.
%         \item If the paper includes experiments, a No answer to this question will not be perceived well by the reviewers: Making the paper reproducible is important, regardless of whether the code and data are provided or not.
%         \item If the contribution is a dataset and/or model, the authors should describe the steps taken to make their results reproducible or verifiable. 
%         \item Depending on the contribution, reproducibility can be accomplished in various ways. For example, if the contribution is a novel architecture, describing the architecture fully might suffice, or if the contribution is a specific model and empirical evaluation, it may be necessary to either make it possible for others to replicate the model with the same dataset, or provide access to the model. In general. releasing code and data is often one good way to accomplish this, but reproducibility can also be provided via detailed instructions for how to replicate the results, access to a hosted model (e.g., in the case of a large language model), releasing of a model checkpoint, or other means that are appropriate to the research performed.
%         \item While NeurIPS does not require releasing code, the conference does require all submissions to provide some reasonable avenue for reproducibility, which may depend on the nature of the contribution. For example
%         \begin{enumerate}
%             \item If the contribution is primarily a new algorithm, the paper should make it clear how to reproduce that algorithm.
%             \item If the contribution is primarily a new model architecture, the paper should describe the architecture clearly and fully.
%             \item If the contribution is a new model (e.g., a large language model), then there should either be a way to access this model for reproducing the results or a way to reproduce the model (e.g., with an open-source dataset or instructions for how to construct the dataset).
%             \item We recognize that reproducibility may be tricky in some cases, in which case authors are welcome to describe the particular way they provide for reproducibility. In the case of closed-source models, it may be that access to the model is limited in some way (e.g., to registered users), but it should be possible for other researchers to have some path to reproducing or verifying the results.
%         \end{enumerate}
%     \end{itemize}


% \item {\bf Open access to data and code}
%     \item[] Question: Does the paper provide open access to the data and code, with sufficient instructions to faithfully reproduce the main experimental results, as described in supplemental material?
%     \item[] Answer: \answerYes{}{} % Replace by \answerYes{}, \answerNo{}, or \answerNA{}.
%     \item[] Justification: We have uploaded anonymized code. 
%     \item[] Guidelines:
%     \begin{itemize}
%         \item The answer NA means that paper does not include experiments requiring code.
%         \item Please see the NeurIPS code and data submission guidelines (\url{https://nips.cc/public/guides/CodeSubmissionPolicy}) for more details.
%         \item While we encourage the release of code and data, we understand that this might not be possible, so “No” is an acceptable answer. Papers cannot be rejected simply for not including code, unless this is central to the contribution (e.g., for a new open-source benchmark).
%         \item The instructions should contain the exact command and environment needed to run to reproduce the results. See the NeurIPS code and data submission guidelines (\url{https://nips.cc/public/guides/CodeSubmissionPolicy}) for more details.
%         \item The authors should provide instructions on data access and preparation, including how to access the raw data, preprocessed data, intermediate data, and generated data, etc.
%         \item The authors should provide scripts to reproduce all experimental results for the new proposed method and baselines. If only a subset of experiments are reproducible, they should state which ones are omitted from the script and why.
%         \item At submission time, to preserve anonymity, the authors should release anonymized versions (if applicable).
%         \item Providing as much information as possible in supplemental material (appended to the paper) is recommended, but including URLs to data and code is permitted.
%     \end{itemize}


% \item {\bf Experimental Setting/Details}
%     \item[] Question: Does the paper specify all the training and test details (e.g., data splits, hyperparameters, how they were chosen, type of optimizer, etc.) necessary to understand the results?
%     \item[] Answer: \answerYes{}{} % Replace by \answerYes{}, \answerNo{}, or \answerNA{}.
%     \item[] Justification: We provided all the details needed to understand the results, including the objective functions, data generation process, hyperparameters, etc. 
%     \item[] Guidelines:
%     \begin{itemize}
%         \item The answer NA means that the paper does not include experiments.
%         \item The experimental setting should be presented in the core of the paper to a level of detail that is necessary to appreciate the results and make sense of them.
%         \item The full details can be provided either with the code, in appendix, or as supplemental material.
%     \end{itemize}

% \item {\bf Experiment Statistical Significance}
%     \item[] Question: Does the paper report error bars suitably and correctly defined or other appropriate information about the statistical significance of the experiments?
%     \item[] Answer: \answerYes{}{} % Replace by \answerYes{}, \answerNo{}, or \answerNA{}.
%     \item[] Justification: We have plotted 95\% and 5\% percentiles among the 100 repetitions. 
%     \item[] Guidelines:
%     \begin{itemize}
%         \item The answer NA means that the paper does not include experiments.
%         \item The authors should answer "Yes" if the results are accompanied by error bars, confidence intervals, or statistical significance tests, at least for the experiments that support the main claims of the paper.
%         \item The factors of variability that the error bars are capturing should be clearly stated (for example, train/test split, initialization, random drawing of some parameter, or overall run with given experimental conditions).
%         \item The method for calculating the error bars should be explained (closed form formula, call to a library function, bootstrap, etc.)
%         \item The assumptions made should be given (e.g., Normally distributed errors).
%         \item It should be clear whether the error bar is the standard deviation or the standard error of the mean.
%         \item It is OK to report 1-sigma error bars, but one should state it. The authors should preferably report a 2-sigma error bar than state that they have a 96\% CI, if the hypothesis of Normality of errors is not verified.
%         \item For asymmetric distributions, the authors should be careful not to show in tables or figures symmetric error bars that would yield results that are out of range (e.g. negative error rates).
%         \item If error bars are reported in tables or plots, The authors should explain in the text how they were calculated and reference the corresponding figures or tables in the text.
%     \end{itemize}

% \item {\bf Experiments Compute Resources}
%     \item[] Question: For each experiment, does the paper provide sufficient information on the computer resources (type of compute workers, memory, time of execution) needed to reproduce the experiments?
%     \item[] Answer: \answerYes{} % Replace by \answerYes{}, \answerNo{}, or \answerNA{}.
%     \item[] Justification: We provide such information in the experiment section. 
%     \item[] Guidelines:
%     \begin{itemize}
%         \item The answer NA means that the paper does not include experiments.
%         \item The paper should indicate the type of compute workers CPU or GPU, internal cluster, or cloud provider, including relevant memory and storage.
%         \item The paper should provide the amount of compute required for each of the individual experimental runs as well as estimate the total compute. 
%         \item The paper should disclose whether the full research project required more compute than the experiments reported in the paper (e.g., preliminary or failed experiments that didn't make it into the paper). 
%     \end{itemize}
    
% \item {\bf Code Of Ethics}
%     \item[] Question: Does the research conducted in the paper conform, in every respect, with the NeurIPS Code of Ethics \url{https://neurips.cc/public/EthicsGuidelines}?
%     \item[] Answer: \answerYes{} % Replace by \answerYes{}, \answerNo{}, or \answerNA{}.
%     \item[] Justification: The research conforms with the NeurIPS Code of Ethics.
%     \item[] Guidelines:
%     \begin{itemize}
%         \item The answer NA means that the authors have not reviewed the NeurIPS Code of Ethics.
%         \item If the authors answer No, they should explain the special circumstances that require a deviation from the Code of Ethics.
%         \item The authors should make sure to preserve anonymity (e.g., if there is a special consideration due to laws or regulations in their jurisdiction).
%     \end{itemize}


% \item {\bf Broader Impacts}
%     \item[] Question: Does the paper discuss both potential positive societal impacts and negative societal impacts of the work performed?
%     \item[] Answer: \answerNA{} % Replace by \answerYes{}, \answerNo{}, or \answerNA{}.
%     \item[] Justification: This is a foundational research paper and there is no direct societal impact of the work.
%     \item[] Guidelines:
%     \begin{itemize}
%         \item The answer NA means that there is no societal impact of the work performed.
%         \item If the authors answer NA or No, they should explain why their work has no societal impact or why the paper does not address societal impact.
%         \item Examples of negative societal impacts include potential malicious or unintended uses (e.g., disinformation, generating fake profiles, surveillance), fairness considerations (e.g., deployment of technologies that could make decisions that unfairly impact specific groups), privacy considerations, and security considerations.
%         \item The conference expects that many papers will be foundational research and not tied to particular applications, let alone deployments. However, if there is a direct path to any negative applications, the authors should point it out. For example, it is legitimate to point out that an improvement in the quality of generative models could be used to generate deepfakes for disinformation. On the other hand, it is not needed to point out that a generic algorithm for optimizing neural networks could enable people to train models that generate Deepfakes faster.
%         \item The authors should consider possible harms that could arise when the technology is being used as intended and functioning correctly, harms that could arise when the technology is being used as intended but gives incorrect results, and harms following from (intentional or unintentional) misuse of the technology.
%         \item If there are negative societal impacts, the authors could also discuss possible mitigation strategies (e.g., gated release of models, providing defenses in addition to attacks, mechanisms for monitoring misuse, mechanisms to monitor how a system learns from feedback over time, improving the efficiency and accessibility of ML).
%     \end{itemize}
    
% \item {\bf Safeguards}
%     \item[] Question: Does the paper describe safeguards that have been put in place for responsible release of data or models that have a high risk for misuse (e.g., pretrained language models, image generators, or scraped datasets)?
%     \item[] Answer: \answerNA{} % Replace by \answerYes{}, \answerNo{}, or \answerNA{}.
%     \item[] Justification: This paper poses no such risks.
%     \item[] Guidelines:
%     \begin{itemize}
%         \item The answer NA means that the paper poses no such risks.
%         \item Released models that have a high risk for misuse or dual-use should be released with necessary safeguards to allow for controlled use of the model, for example by requiring that users adhere to usage guidelines or restrictions to access the model or implementing safety filters. 
%         \item Datasets that have been scraped from the Internet could pose safety risks. The authors should describe how they avoided releasing unsafe images.
%         \item We recognize that providing effective safeguards is challenging, and many papers do not require this, but we encourage authors to take this into account and make a best faith effort.
%     \end{itemize}

% \item {\bf Licenses for existing assets}
%     \item[] Question: Are the creators or original owners of assets (e.g., code, data, models), used in the paper, properly credited and are the license and terms of use explicitly mentioned and properly respected?
%     \item[] Answer: \answerNo{} % Replace by \answerYes{}, \answerNo{}, or \answerNA{}.
%     \item[] Justification: The code for this paper is adapted from the code for \citep{chen2023generalized}, which does not have license. We have obtained the permission from the authors of \citep{chen2023generalized} and cited this paper in the experiment section. 
%     \item[] Guidelines:
%     \begin{itemize}
%         \item The answer NA means that the paper does not use existing assets.
%         \item The authors should cite the original paper that produced the code package or dataset.
%         \item The authors should state which version of the asset is used and, if possible, include a URL.
%         \item The name of the license (e.g., CC-BY 4.0) should be included for each asset.
%         \item For scraped data from a particular source (e.g., website), the copyright and terms of service of that source should be provided.
%         \item If assets are released, the license, copyright information, and terms of use in the package should be provided. For popular datasets, \url{paperswithcode.com/datasets} has curated licenses for some datasets. Their licensing guide can help determine the license of a dataset.
%         \item For existing datasets that are re-packaged, both the original license and the license of the derived asset (if it has changed) should be provided.
%         \item If this information is not available online, the authors are encouraged to reach out to the asset's creators.
%     \end{itemize}

% \item {\bf New Assets}
%     \item[] Question: Are new assets introduced in the paper well documented and is the documentation provided alongside the assets?
%     \item[] Answer: \answerNo{} % Replace by \answerYes{}, \answerNo{}, or \answerNA{}.
%     \item[] Justification: The paper does not release new assets.
%     \item[] Guidelines:
%     \begin{itemize}
%         \item The answer NA means that the paper does not release new assets.
%         \item Researchers should communicate the details of the dataset/code/model as part of their submissions via structured templates. This includes details about training, license, limitations, etc. 
%         \item The paper should discuss whether and how consent was obtained from people whose asset is used.
%         \item At submission time, remember to anonymize your assets (if applicable). You can either create an anonymized URL or include an anonymized zip file.
%     \end{itemize}

% \item {\bf Crowdsourcing and Research with Human Subjects}
%     \item[] Question: For crowdsourcing experiments and research with human subjects, does the paper include the full text of instructions given to participants and screenshots, if applicable, as well as details about compensation (if any)? 
%     \item[] Answer: \answerNA{} % Replace by \answerYes{}, \answerNo{}, or \answerNA{}.
%     \item[] Justification: The paper does not involve crowdsourcing nor research with human subjects.
%     \item[] Guidelines:
%     \begin{itemize}
%         \item The answer NA means that the paper does not involve crowdsourcing nor research with human subjects.
%         \item Including this information in the supplemental material is fine, but if the main contribution of the paper involves human subjects, then as much detail as possible should be included in the main paper. 
%         \item According to the NeurIPS Code of Ethics, workers involved in data collection, curation, or other labor should be paid at least the minimum wage in the country of the data collector. 
%     \end{itemize}

% \item {\bf Institutional Review Board (IRB) Approvals or Equivalent for Research with Human Subjects}
%     \item[] Question: Does the paper describe potential risks incurred by study participants, whether such risks were disclosed to the subjects, and whether Institutional Review Board (IRB) approvals (or an equivalent approval/review based on the requirements of your country or institution) were obtained?
%     \item[] Answer: \answerNA{} % Replace by \answerYes{}, \answerNo{}, or \answerNA{}.
%     \item[] Justification: The paper does not involve crowdsourcing nor research with human subjects.
%     \item[] Guidelines:
%     \begin{itemize}
%         \item The answer NA means that the paper does not involve crowdsourcing nor research with human subjects.
%         \item Depending on the country in which research is conducted, IRB approval (or equivalent) may be required for any human subjects research. If you obtained IRB approval, you should clearly state this in the paper. 
%         \item We recognize that the procedures for this may vary significantly between institutions and locations, and we expect authors to adhere to the NeurIPS Code of Ethics and the guidelines for their institution. 
%         \item For initial submissions, do not include any information that would break anonymity (if applicable), such as the institution conducting the review.
%     \end{itemize}

% \end{enumerate}


\end{document}
